\documentclass[11pt,openany]{book}

% Idioma y codificación
\usepackage[spanish]{babel}
\usepackage[T1]{fontenc}
\usepackage[utf8]{inputenc}

% Página
\usepackage[
  paperwidth=6in,
  paperheight=9in,
  inner=0.9in,
  outer=0.75in,
  top=0.8in,
  bottom=0.8in
]{geometry}

% Tipografía y párrafos
\usepackage{ebgaramond}
\usepackage{microtype}
\usepackage{setspace}
\usepackage{emptypage}
\usepackage{tikz}
\usepackage{pgfornament}
\onehalfspacing

\setlength{\parindent}{1.5em}
\setlength{\parskip}{0.5em}

% Encabezados
\usepackage{fancyhdr}
\pagestyle{fancy}
\fancyhf{} % limpia encabezados y pies
\fancyhead[LE]{\thepage}
\fancyhead[RO]{\thepage}
\renewcommand{\headrulewidth}{0pt}
\renewcommand{\headrulewidth}{0pt}
\setlength{\headheight}{13.59999pt}

\usepackage{titlesec}

\titleformat{\chapter}[hang]
  {\Large\bfseries\centering}
  {\Roman{chapter}.}
  {0.5em}
  {}

\titlespacing*{\chapter}
  {0pt}
  {2cm}
  {1cm}

% Primera página de capítulo sin encabezado
\makeatletter
\let\ps@plain\ps@empty
\makeatother

% Metadatos
\title{\Huge Demens et Infamis}
\author{Rodrigo Blanco Betes}
\date{}

\begin{document}

% Portada
\frontmatter
\begin{titlepage}
  \centering
  \vspace*{4cm}
  {\Huge\scshape Demens et Infamis\par}
  \vspace{1.5cm}
  {\Large Rodrigo Blanco Betes\par}
  \vfill
\end{titlepage}

% Dedicatoria
\cleardoublepage
\thispagestyle{empty}
\vspace*{6cm}
\begin{flushright}
  \textit{¿Cómo puede la belleza ser algo tan espantoso, y al mismo tiempo lleno de misterio? Es la eterna lucha de Dios contra el Demonio, cuyo campo de batalla es el corazón del hombre.}

  F. M. Dostoievski 

\hfill

  \textit{El hombre es una cuerda finísima entre el superhombre y el animal: una cuerda encima del abismo\dots}

  F.Nietzsche 

\end{flushright}
\newpage

% Índice opcional
% \tableofcontents

\section*{Nota del autor:}

El presente escrito no es una obra sino un ejercicio. No se encontrarán pues aquí el acabado ni las facilidades que se presumen en un escrito destinado y pensado para sus lectores; este libro se ha escrito exclusivamente para su autor, y exclusivamente porque así se lo ha impuesto el estado actual de su pensamiento.  

Queda advertido así el lector de que lo que aquí hubiere de encontrarse en ningún caso tendría por qué resultarle de su agrado, y de que, si así y todo quisiera adentrarse en tal lectura, él solo responde de la pérdida de tiempo y de los otros posibles efectos negativos que ésta pudiera ocasionarle. No encontrará una trama coherente, intrigante ni bien pulida, ni recursos literarios grandilocuentes, ni si quiera son atractivos, en mi opinión, los nombres de los personajes; si se diera alguno de estos casos se trataría de un accidente, ajeno a la intención del autor. La redacción le resultará por momentos áspera y difícil; de hecho creo que este es un libro para leerlo despacio, si acaso se merece que alguien lo lea. El hecho de que el contexto del libro es el mundo romano, es sólo algo accesorio, inevitable. Simplemente este es el contexto que se ha figurado. Esto no es tampoco una novela histórica; no pretende situarse en ningún momento concreto de la época julio-claudia, flavia, severa ni nada del estilo. Lo único que hace la novela propiamente romana es su punto de decadencia; debería por tanto tratarse de algún momento de la decadencia del Imperio, pero eso es vago e impreciso. No debe venirse aquí, en suma, en busca de precisión histórica. 

Advertidos los aspectos negativos del escrito, y eximido por ello de culpa su autor, puedo pasar a un apunte de lo que sí se podrá esperar del mismo. \textit{Demens} se refiere al comportamiento de los personajes, chocante y extravagante de continuo. Hay un punto de desorden, algo de inesperado y absurdo, en cada pasaje del libro. Todo esto contrasta con el narrador, que viene a ser la otra cara del espejo. El cuestor Cayo Flaminio es la representación, si bien exagerada y caricaturesca, del orden romano, que no puede estar más alejado del espíritu estrafalario de la mayoría de los demás personajes. Es narrador y al mismo tiempo juez de la obra hasta su desenlace, lo \textit{infamis}: la desgracia que sucede al absurdo. Semejante tarea debería encomendarse seguramente a un gran hombre, lo cual no es el caso; el cuestor es corto de miras, maniático y obtuso. A pesar de ello cumple su función, ya que en todo caso, y contra lo que pudiera parecer, es un gran observador; ésta es la única virtud a la que en mi opinión está obligado un buen narrador. Su pobreza es seguramente de espíritu, pero no de inteligencia. 

Inteligencia y profundidad de espíritu sí que convergen en el protagonista de la obra: Marco Valerio Cátulo. Extravagante, impulsivo, si no malvado por momentos, extraño y distante, ¿por qué es mi héroe, por qué considerarlo sublime? Pues, y líbreme Dios de las malas lenguas, porque yo juro que no apruebo casi ninguna de las barbaridades que cometerá este individuo en mi obra, Marco Valerio me parece un personaje admirable, del máximo interés. Su primera característica esencial es su desvinculación moral con la norma socialmente establecida. Cuando se pierde el cómodo amparo del rebaño, aparecen el miedo, la verdadera soledad, la angustiosa y genuina duda existencial, pero se despeja el único camino posible a la verdad y a la auténtica felicidad. Cuando uno se desenvuelve en este camino se descubren tanto su grandeza como su villanía, desprovistas de la trampa de la norma. ¿Cuál de estas dos atesora Marco Valerio? Eso no se decidirá en ningún momento del libro. El cuestor representa precisamente la norma social, para la cual todo lo que se sale de ella es infame, resulte como resulte. Yo no juzgaría con la tosquedad del rebaño, sin embargo, a alguien tan elevado y plagado de importantes intuiciones, así como de los inevitables desvaríos y contradicciones. Ése es el cometido de la obra, seguramente: el de juzgar, únicamente ante el tribunal de la moral, las absurdas infamias de Marco Valerio, con todos sus matices tan diversos y problemáticos. Estoy seguro de que el lector será un magistrado mucho más sofisticado que el infeliz y pobre narrador. 

No debemos adelantarnos sin embargo; todo se dará cuando llegue el momento. Valgan estos apuntes simplemente para que, como antes ya he dicho, uno sea consciente de lo que se va a encontrar en el libro antes de decidir si quiere empezar a leerlo. Por lo demás, estas explicaciones preliminares son sin duda innecesarias; si las he incluido, probablemente ha sido para tener un modo de decir que el libro que he escrito no es bueno. Por tanto quizá este prefacio sea lo más aceptable del texto, por contener al menos ese cierto punto de sinceridad artística, aunque esto tampoco lo dota en ningún caso de calidad; se trata, sencillamente, de un ejercicio de falsa modestia. Así que pongámonos manos a la obra, ya que, no habiendo sido esta nota capaz de salvar el libro, igual al menos sí es el libro capaz de justificar la nota.
 

% Cuerpo
\mainmatter
\small
\textit{Al ilustrísimo cónsul Flaminio,}

\textit{Escribo, hermano, en los más graves términos que pueden concebirse. Corresponde a su alteza imperial durante las presentes fechas, como ya sabrás, la recepción del informe del procónsul de Grecia sobre el estado y la gestión de la región. Éste no será sin embargo recibido por su majestad, ni enviado en modo alguno, y en su lugar te hago llegar la espantosa noticia de lo sucedido durante el ejercicio; has de saber que su triste protagonista es nadie menos que Marco Valerio Cátulo, el hijo de su alteza y heredero de su trono, al menos aún sobre el papel. Tu labor será la más ingrata que queda en estos terribles casos; habrás de elegir, entre cuanto se te haga saber por la presente, las palabras adecuadas para transmitirle su contenido al emperador, quien desde ahora doy fe de que será prisionero en una jaula de oro, si acaso sus fuerzas le permiten seguir en ella. Confío en que tu conocimiento y tu cercanía al desdichado son lo que requerimos para minimizar el inevitable daño causado, tanto a él como a nosotros los políticos y a todo el resto del mundo romano.}

\textit{He descrito de forma realmente minuciosa todos los detalles que puedan tener relación con el suceso, sencillamente para facilitar tu tarea, haciéndote disponer de cuánta información estimares oportuna. De todos modos, en ningún caso te dirijas a su majestad sin el completo conocimiento de todo lo sucedido; en un principio, podrá parecer que exagero ante un simple escándalo, pero los hechos han derivado inesperadamente en el ominoso resultado final. ¿Cómo se habrán dado las circunstancias para que ahora mismo, a todos los efectos, sea yo el procónsul, y lo que tenga que anunciar consista en una desgracia descomunal, que excede con mucho cuanto habríamos podido prever o imaginar? El gobernador ha muerto y su desgraciado verdugo ha desaparecido. No me adelantaré sin embargo a la narración de los hechos. Quizá otra parte del cometido que injustamente te he asignado es la paciencia con la que te pido que recibas relación de todos ellos. Sin embargo, escribiendo como hermano a la vez que como magistrado, al cargo de uno más de los eslabones que organizan el mundo conocido, he de reconocer que, si bien desgraciados, los hechos que te haré saber son extraordinarios, y que bien empleados estarían la paciencia y el tiempo que pudieras dedicarles.}

\textit{Es tu humilde servidor, }

\textit{El cuestor Cayo Flaminio}

\chapter*{Prólogo. A seis días de calendas de marzo}

Me reprochará un hombre de tan cuantiosas ocupaciones que en el principio de mi relato no aparezca su protagonista, ni tampoco suceda nada con relación a lo que narraré posteriormente. Es, sin embargo, a mi juicio, tan relevante este inicio de mi escrito como cuanto diré en cualquier otro momento; la rueda de la fortuna se conforma de rígidos engranajes, de forma que si uno se hunde en las tinieblas, arrastra con él a sus contiguos y a sus semejantes. Así, el día en que el loco ambulante trajo su lamentable falta de cordura a la imperial Atenas, la dejó prendida de ella sin remedio, y ni notarlo todo esto en su momento, como me correspondía desde mi puesto de cuestor, ni las medidas que tomé al respecto, sirvieron para revertir la situación. 

Era un día sorprendentemente cálido para la época del año, febrero, en la emblemática colina del Areópago; a falta de juicios, asambleas ni eventos de tal índole, sólo se veía a la ciudad de Atenas salida a sus calles, disfrutando del clima primaveral. Yo mismo me encontraba, casualmente y sin que ahora mismo recuerde el motivo, en el lugar. Fue entonces cuando apareció él; ¡condenado sea por la absoluta eternidad! Este tipo de locos son uno de los productos excepcionales de la caprichosa naturaleza; un día el libro equivocado llega a sus manos, en el que leen absurdas historias, por lo general judías, de profetas que escuchan de primera mano los designios de los dioses, al siguiente creen que sería excitante vivir como alguno de esos profetas, y finalmente se despiertan un día predicando en nombre de tal o cual divinidad contra los maliciosos vicios de los mortales y los peligros que éstos acarrean, olvidando el que ellos mismos corren de hacer el ridículo. Obviamente, después de lo sucedido, lo detuve y lo interrogué, pero como era de esperar de su versión de los hechos no obtuve nada verosímil ni claro, luego me quedaré para siempre sin saberlo. Debía de haber andado al menos una semana desde su hogar, ya que al presentarse en la ciudad su aspecto era deleznable; vestía unos sucios harapos, estaba cubierto de arañazos y quemaduras, descalzo, despeinado y greñoso, y se movía lentamente y con dificultad. No estaba tan fatigado, sin embargo, como para no poder subirse, ante la inexplicable pasividad de los soldados allí enviados por el prefecto, al estrado del Areópago. Allí, como los oradores de la primera democracia, no sé si creyéndose que era Sócrates, Pericles o Alcibíades, empezó a berrear: 

“¡Absurdos sois, absurdos!... ¡Absurdos somos todos; pero vosotros sois peores y más absurdos que yo!” 

Los paseantes, llenos de desconcierto y no encontrando, como es natural, otro entretenimiento, tornaron todos sus cabezas hacia el estrado. Allí el pobre loco hizo varias extrañas gesticulaciones y sonidos, de los que sin encontrarles interpretación razonable no daré más relación, contento quizá de ver cómo su labor mesiánica se veía recompensada por la atención de la masa, ante cuyo sepulcral silencio continuó con sus extraños gritos: 

“¡Los dioses os odian, y se ríen de vosotros!... También de mí\dots ¡pero más aún de vosotros! Ellos son fuertes e inteligentes, vosotros frágiles y manipulables, ellos entienden la vida y la humanidad, y vosotros sois caricaturas de personas\dots ¡y vuestras vidas son ridículas caricaturas de la misma existencia: la existencia que nunca entenderéis!” 

El llanto de un niño interrumpió los gritos de este pobre loco. La madre del niño, avergonzada de haber llamado la atención, abandonó el lugar rápidamente con él en brazos.  

El loco se calló un momento. Encorvado, apoyando los brazos sobre la barandilla del estrado, miró primero hacia el cielo despejado, y luego a varios entre la multitud, poniendo de forma realmente extraña sus ojos como platos. La expectación por ver lo que iba a decir se había acrecentado; nadie se movió de su sitio ni dijo palabra, hasta que él, gesticulando de nuevo con su boca, dijo: 

“Sois espantosos y ridículos, porque formáis parte de la masa\dots la masa que de todo se apropia; la masa que pierde el talento del hombre superior, y que esconde y disimula la torpeza del simple; la masa espantosa\dots pero que no sabe nada de espanto\dots ¡Está acostumbrada a sufrir, pero no entiende lo que es el sufrimiento!” 

Se oyó un solitario grito de alabanza desde el público, pero el loco casi no le hizo caso; con la mirada aún más perdida, gesticulando y gruñendo mientras hablaba, dijo: 

“La masa\dots la masa también\dots —entonces paró un momento de hablar porque no pudo eludir un repentino ataque de risa— la masa también tiene sus héroes\dots —ahora calló al suelo de la risa— ¡y son ridículos!” 

Salió entonces del estrado que ocupaba, y empezó a recorrer la colina totalmente errático. Por donde pasaba le abrían paso las gentes, que habían llegado a ocupar casi la totalidad del lugar. No sé si se le apartaban por respeto a su figura o, lo que es más probable, por el hedor que emanaba. Parándose ante un hombre, que llevaba la túnica de los maestros de escuela, empezó a chillarle, provocando que al de poco escapara del lugar espantado: 

“¡Santa santísima locura! ¡Qué hacer contra la locura y el espanto\dots qué hacer! Sois ridículos porque sois pobres de mente y porque os contentáis de hacer lo que veis hacer\dots ¡¡liberaos!! Sed libres, escuchaos en vuestra verdad\dots¡sed absurdos y libres, mejor que absurdos y encadenados! —lo interrumpió entonces, por un momento, si mal no recuerdo, un breve ataque colectivo de tos, cuando yo ya me acercaba, errante entre las masas, a la patrulla del prefecto— ¡¡Estáis dormidos, hijos míos, pero vengo yo para despertaros; estáis atados, pero yo he venido a desencadenaros, alabadme y encontraréis vuestra salvación!!” 

Fue entonces cuando unos jóvenes, auténticos granujas, se acercaron, en un grupo de cuatro o cinco, a la, él desde luego que lo era, soberana caricatura de profeta o de lo que diablos pretendiera ser el hombre; dijo entonces uno de éstos, que probablemente se habrían escapado del parvulario, o los habían expulsado por sus malas prácticas: 

“¡Eh, pueblo de Atenas! ¿Habéis escuchado lo mismo que nosotros de la boca de este loco?... ¡Seamos libres! ¡¡Libres!! ¡Este enviado del Olimpo nos promete la salvación si gozamos de nuestra libertad como es debido!” 

Entonces, entre los cinco que iban, empezaron a sacar piedras, que llevaban en bolsas escondidas, y a freír a pedradas al pobre chalado, que esquivaba los golpes como buenamente pudo. Esto desató risas y un lamentable griterío entre la multitud; había cundido definitivamente el desorden. Fue entonces cuando, por suerte, llegué ante los hombres del prefecto; me presenté ante ellos como cuestor, diciéndoles que aquél era un escándalo inaceptable y les exigí que represaliaran a los responsables. La intervención fue rápida y eficiente y los jóvenes fueron reducidos al instante; obviamente, ordené que fueran enviados al calabozo. Seguramente eran hijos de familias ricas pero laxas y despreocupadas, lo cual no me molesté en investigar, y supuse que una semana de presidio sería el castigo adecuado para su comportamiento. La multitud fue disuelta, sin que se requiriera mayor uso de la fuerza, y entonces sólo quedaba por decidir qué hacer con el pobre diablo que había originado la situación. Que fue detenido y acusado de tumulto y de escándalo público no hace falta ni que lo diga; sin embargo, no pensé que un año en las galeras, o tres en las minas de Córdoba, como suele hacerse en estos casos, fueran lo más adecuado para un loco, además de decrépito, y fue así que decidí mantenerlo encerrado para ocuparme yo mismo del caso. Hablaré más adelante, en todo caso, de cuanto dio de sí el sumario de este curiosísimo caso. 

En lo que a nosotros concierne, debo decir que, como ya adelanté, mi intervención tardía fue un fracaso; al día siguiente, fuera con sorna o, por inverosímil que parezca, admiración, toda Atenas hablaba de la aparición del loco. Habíamos sido invadidos por la niebla de la locura, que se extendió como la peste por la ciudad, y acabaría por oscurecer el brillo de nuestra más preciada joya.

\chapter{A un día de calendas de marzo}

Los mejores brillantes están hechos de carne y hueso, repugnantes de por sí, si no fuera por el inexplicable sistema que forma el organismo; a Marco Valerio no lo has visto, si mal no recuerdo, desde que a sus cuatro años el César decidió enviarlo lejos de Roma, de lo cual hace ya veinte. ¿Qué mejor que la Atenas de Fidias, Sócrates y Alcibíades, cuna de la democracia, a la que lo suyo debe el equilibrio político romano; qué mejor lugar que la capital de la política, el arte y el pensamiento, para forjar el carácter del que está llamado a gobernar el mundo civilizado, culto y avanzado? ¿Qué mejor destino que el corazón de la Grecia conquistada que conquistó al bárbaro conquistador? Yo no me opuse a tal decisión entonces, ni encontré con el tiempo motivos para hacerlo; el niño creció, rodeado de lo mejor de la academia griega, y mostró un progreso asombroso, hasta convertirse en un joven brillante, conocedor eminente del arte del buen gobierno, la filosofía y la historia: todo lo que un césar ha de saber. 

Ahora bien, todo lo que de su inteligencia era propio de un buen gobernante, dejaba de serlo en su extraño carácter, rasgo por desgracia frecuente entre los emperadores de Roma. Los galenos habrían dicho del chico que era hipersensible, quiera decir eso lo que quiera decir. por poner un ejemplo: la melodía que a uno puede sencillamente agradarle, a él le llegaba del todo a las entrañas, del mismo modo que se alegraba, entristecía o esperanzaba muchas veces sin motivo justificado, a ojos de los que le rodeábamos. Al mismo tiempo era un muchacho ciertamente distante; no era raro encontrarlo en silencio, caviloso, caminando sin rumbo con la cabeza gacha, aún más en su temprana juventud, sobre la docena de años. La forma de la que yo lo veía me llamaba profundamente la atención, sobre todo en contraste con los rudos y poco comedidos jóvenes de esas edades. Tal contraste de caracteres podría haberle jugado una mala pasada a Marco Valerio en la escuela, a la que acudía además de las lecciones que recibía de sus maestros particulares. La crueldad de los niños es implacable, y en otros casos yo he conocido que este tipo de formas de ser eran tildadas de tristes y afeminadas; ¡cuánto me he tenido que enfrentar yo como magistrado a estas malvadas costumbres de los niños, con las consiguientes quejas de sus padres! Sin embargo, es curioso en este caso cómo Marco Valerio era realmente querido por los demás jóvenes; incluso diría que comprendían y admiraban su carácter. Además, guardábamos en secreto el que se trataba del hijo del emperador, por lo cual esta afección de los demás hacia él no era interesada, sino auténtica. En el fondo yo creo que, a pesar de sus rarezas, el chico sabía comportarse ante los otros, y los conocía profundamente, y además siempre que hablaba, independientemente de con quién lo hiciera y del tema, demostraba cuidado e inteligencia. Además, y quizá con lo que verás que diga más adelante piensas que he perdido el juicio pronunciando estas palabras; además, decía, creo que el chico era en el fondo de buen corazón, y no es raro que, incluso con rarezas, alguien de buen talante reciba un buen trato. Yo creo que lo habitual es que sea apreciado, aunque no siempre tomado en serio. 

Si me dices que es extraño, en Marco Valerio, que alguien de su cariz extravagante tuviera buen trato con sus colegas en la escuela, más aún lo tiene que ser su amistad con Lucio Junio Bruto. Se trataba nada menos que del hijo de Cneo Junio Bruto y Julia Severa, a quienes habrás de conocer por su estancia en Neápolis justo antes de su definitivo traslado a Grecia. Recordarás los ataques de furia de Cneo Junio, los cuales son igual de poco extraños, pero menos graves, en su hijo; diría que éste es más bien brusco que furioso. Lo lamento por la pobre Julia Severa, una gota de dulzura en la marea de barbaridad que es esa familia. Lo que venía diciendo, es que el bruto y frívolo Lucio Junio, por absurdo que parezca, y el sensible y caviloso Marco Valerio, de edades similares, se conocieron hará unos cinco años, para acabar siendo amigos inseparables. La verdad yo diría que ambos son la exacta antítesis el uno del otro, y si nunca han acabado a gritos y mamporros, es porque hasta tal punto llegan sus diferencias, que nunca podrían hacer ni pensar ni sentir el uno y el otro lo mismo en el mismo momento, aunque fuera para odiarse y discutir.

Marco Valerio fue un joven responsable y respetuoso con las normas establecidas hasta el inicio de su relación con Lucio Junio. Las trastadas de éste en su niñez se relacionaban con bromas pesadas, novatadas y gansadas varias, pero en su juventud ya lo hacían con la bebida, las tabernas y los bajos fondos. Los patricios romanos no suelen dejarse llevar por estas malas costumbres, y por tanto estos dos recurrían a una táctica peculiar, a la vez terriblemente peligrosa, para pasar desapercibidos; se vestían con ropas de plebeyos griegos, escondiendo entre ellas una daga para el caso de que se vieran inmiscuidos en una situación violenta. Entonces cabalgaban, a lo largo del Muro Largo, hasta el barrio industrial del Pireo; como no conoces Atenas, debo decirte que este lugar está lleno de tascas mugrientas y antros varios de este estilo, donde uno se emborracha hasta la saciedad, se deja el dinero en juegos estúpidos y se acaba matando a cuchilladas con otros indeseables. Según lo que me dijo más tarde Lucio Junio, no iban allí a participar en estas salvajadas, sino simplemente a “divertirse un poco” y  “disfrutar del espectáculo” que todo esto debía de producirles. Fue en estas nonas de marzo que Marco Valerio debió de quejarse de la monotonía de tantos días con el mismo tipo de burdas diversiones; entonces Lucio Junio le dijo: 

“¡Será posible, Cátulo, lo poco que tardas en cansarte de cualquier cosa! Sea pues, tenía una sorpresa reservada; pensé que para un día más interesante, pero tendrá que ser para hoy.” 

Entonces condujo a Marco Valerio por una serie de callejones, orientado como si hubiera diseñado él mismo aquellas barriadas del demonio. Al final, comprobando que nadie les hubiera seguido ni les estuviera viendo, se coló, pidiendo a su camarada que lo siguiera, por un pequeño túnel, disimulado dentro la pared de ladrillo de una ínsula; lo cruzaron y se encontraron en un paraje solitario y ciertamente hermoso al borde del mar. Las ínsulas del barrio se agolpaban contra la costa, dejando un pequeño camino de roca gris entre ellas y el mar, al fondo del cual se veían todos los barcos que entraban y salían en los muelles del Pireo: 

“¿Vas a quejarte —dijo Lucio Junio— también de esto? No sería raro que encontrara, tu paladar exquisito, un sabor desagradable incluso en este paraje tan hermoso y relajante.” 

Marco Valerio miró con atención en todas direcciones, hasta que dijo: 

“La verdad, Lucio, es que has encontrado por fin un lugar interesante, en vez de uno lleno de estúpidos y borrachos. ¿Podríamos acercarnos en esta dirección—dijo Marco Valerio señalando hacia un lado—, viendo más de cerca los muelles?” 

“Desde luego; de todos modos, ¿no te interesa saber cómo he encontrado este lugar?” 

“Qué más me da.”

“¡Ah...!” 

Lucio Junio accedió, sin embargo, a la petición de Marco Valerio, y ambos avanzaron por el camino; entre observaciones de este último sobre la fauna del lugar y quejas del otro por los inevitables golpes de las olas del mar, en un camino en la misma costa, llegaron a su final. Se veía perfectamente cómo entraban las galeras: no sé por qué, pero Lucio Junio conocía perfectamente el origen y la misión de cada una de ellas, fueran comerciales o militares, y sorprendió a Marco Valerio con semejantes conocimientos. Donde estaban había un pequeño mirador, desde el que se veía, a apenas unos diez metros, el más septentrional de los muelles de Atenas: 

“¿A quién pertenece este muelle? —preguntó Marco Valerio apoyado en la barandilla del mirador—.” 

“A nadie menos que Avidio Régulo. Como estás tan al tanto de la actualidad, me imagino que habrás oído de quién se trata\dots” 

“Ni palabra.” 

“Avidio Régulo es el mercader más rico y odiado del mundo romano. Que los mercaderes piensan sólo en el oro, el cual es su única dicha posible, que incumplen cualquier código de conducta para conseguir su paupérrimo objetivo de la riqueza, todo eso es cosa sabida. Sin embargo, en cuanto a su ambición, sus repugnantes triquiñuelas, su mal trato hacia los esclavos, sus robos del tesoro del estado, disimulados con el porcentaje adecuado en favor del funcionario correspondiente, como no podía ser de otra manera; en cuanto a tantas otras cosas, nadie destaca más que Avidio Régulo. Que es rico hasta la saciedad, no hace falta que lo diga; pues si no lo fuera siendo tan malvado, entonces tendría que ser tonto de remate, en cuyo caso no habría ningún motivo para que fuera famoso.” 

Marco Valerio observó detenidamente el lugar, donde había esclavos cargando fardos y ánforas en un barco: 

“Bueno, éstos no parecen tan maltratados; otras veces los he visto con más marcas en la piel.” 

“En serio te digo —respondió riéndose Lucio Junio— que no sabes sobre lo que haces bromas. Si éstos no tienen marcas será porque el que se dedicaba a darles latigazos habrá muerto de extenuación. He oído que el año pasado en Capua hubo un levantamiento en sus villas; sus esclavos, incluso teniendo el famoso precedente de Espartaco, se hartaron de sus malas prácticas y se rebelaron. Al de cinco días habían tomado la finca, incluso secuestrado a dos sobrinos de Régulo como rehenes; éste llegó de inmediato, escoltado por una cohorte de legionarios. Los esclavos mataron a los rehenes, nada más ver que se encontraba él en el lugar, y le lanzaron sus cabezas untadas en melaza; me imagino que se vieron perdidos, y que sencillamente, la derrota, quisieron consolarla el agridulce fruto de la venganza. De inmediato los legionarios asaltaron la finca; bien, que todos los que trabajaban en ella acabaron muertos, o torturados hasta la muerte, y del mismo modo sus mujeres y sus niños; lo curioso es cómo después disecó sus cuerpos y los repartió por su terreno. De este modo, los luego volvieron a trabajar esa tierra, veían cada día como el espantapájaros de tal viña era un hombre disecado, como en un merendero había colocada una familia quemada viva, como si quisiera simular que estaban descansando. He oído también que con algunos cadáveres no sabía ya que hacer, y entonces los hizo pedacitos y ordenó que fueran utilizados de abono para los huertos; ¡imagínate esto! No pienses en cualquier caso que su impiedad se limita a los esclavos; hace cinco años produjo una hambruna especulando sobre el precio del trigo, con lo que se hizo de oro. Sin embargo al César no debió de gustarle esto, ni a la ciudadanía romana en general, porque entonces comenzaron a controlar el precio del trigo, y él tuvo que exiliarse en Grecia, por la vergüenza y la pérdida de amistades. ¡Vaya exilio de oro, después del expolio que había hecho en Roma, venir a gastarlo todo en Grecia!” 

Ciertamente, no sé si Lucio Junio exageraba, y todo esto eran simples rumores, o la avaricia de Avidio Régulo es tan espantosa como se la narró entonces a Marco Valerio. En cualquier caso, creo que este mercader sin entrañas es un peligro para la paz y la prosperidad de Roma; las páginas siguientes no lo mostrarán ajeno a esta historia, y de hecho yo mismo acabé por interrogarlo mientras investigué el caso. Pragmáticamente no tomé mayores medidas contra él, lo cual sin embargo merecerá una pronta revisión.

Marco Valerio se quedó observando el muelle, ciertamente interesado de lo que le había contado Lucio Junio, que sonrío contento de haber podido demostrar sus buenas informaciones, incluso entreteniendo a su exigente compañero. Éste lanzó una mirada más al muelle, antes de acercarse a Lucio Junio: 

“¿Se te puede preguntar una cosa?” 

“¿Cuál?” 

“Quiero decir si podré preguntártela sin que te alteres. Me refiero a esa jovencita con la que hablabas el otro día tan alegremente, claro.” 

“Ah\dots ¡es eso! ¿Y qué quieres saber de ella, si se puede saber?” 

“Bueno —Marco Valerio dejó pasar un tiempo antes de responder, mirando con cierta sorna a Lucio Junio—, pues podría interesarme como se llama, de dónde viene, y qué puedes querer de ella. ¿Es una eminencia de la comedia, y quieres que te enseñe a actuar en el teatro?” 

“¡Ah, eres insufrible! ¿Qué te importará? ¡Siempre igual con el señorito Valerio Cátulo!” 

Entonces se puso a andar malhumorado de un lado a otro, hasta volver donde estaba y responder: 

“Está bien, te diré todo, si prometes no decir nada tú. Ella —fue respondido por un asentimiento de Marco Valerio— se llama Albina, y no te voy a decir para qué he andado con ella\dots ¡será posible que preguntes así estas cosas!” 

“¿Ella te quiere?” 

Lucio Junio reaccionó con un suspiro: 

“Sin duda; ¿cómo no iba a quererme a mí?” 

“Pero no me lo has dicho todo, y espero que no estés tratando de ocultármelo; ¿se llama Albina y nada más? ¿Me estás diciendo que no es de buena familia?” 

“Ah\dots ¡Cátulo, siempre igual de difícil!”

La furia de Lucio Junio provocó una sonrisa en Marco Valerio:

“¡Ajá! A ver si adivino\dots ¿familia de libertos?”

“Oye, ¿por quién me tomas? Estaré loco, pero no soy idiota; creo que es équite.”

“¡Équite! ¿Quieres decir que pretendes que el implacable Cneo Junio y la fuertemente tradicional Julia Severa bendecirán tu unión con una familia équite? ¿O pretendes verla a escondidas en el campo por el verano?”

“Bah\dots ¡me ocuparé yo de todo eso!”

“No me digas que te ocuparás de eso; perderás la fortuna y el favor de tu familia, tus amigos te despreciarán, se hablará de ti y de tu caso en la calle, se reirán de ti hasta los plebeyos: se hablará de Junio Bruto, el tonto que se enamoró y casó con una équite de mala manera, perdió su fortuna, y entonces ella le abandonó por otro tonto con dinero, o por un caballero quizá sin dinero pero al menos no un rematado idiota.”

“¡Será posible que seas tan exagerado! Además, ¿quién te ha dicho que vaya a casarme con ella? ¿Casarme con ella, tú qué te piensas? ¿Y qué tontería es esa de que me va a abandonar?” 

“¡Por favor; piensa un poco! ¿Te crees que ella te quiere por tu cara bonita y tus formas encantadoras, por tu elegancia y tus gustos refinados? Querrá que te cases con ella para llevarse la mitad del dinero, y si no te casas, ¡pues se irá con cualquier otro! ¡La vi con mis propios ojos, y era capaz de notarlo a leguas! ¿Es que hablasteis de vuestras vidas, de la inmensidad del mundo, de la alegría del baile y la música? ¡Seguro que no, y que te preguntó quiénes eran tus padres y cómo de ricos y poderosos eran, y tú que eres tan impertinente, le dijiste todo!” 

Lucio Junio estaba fuera de sí, y comenzó a moverse en todas direcciones, tirando piedras y gesticulando, mientras Marco Valerio se tumbó en una roca con una sonrisa de oreja a oreja. ¿No es increíble, cómo entonces este chico era una maravillosa muestra de sensatez y respeto por las buenas prácticas? Lucio Junio se dio la vuelta hacia él entonces: 

“¡Eres insufrible Cátulo! ¡¡Insoportable del todo!!” 

“¡Desde luego que soy insoportable! Soy muy malo, y muy malpensado, así que no me figuro nada bueno de esa Albina\dots ¡mala familia, y seguro que malas costumbres también!” 

Esto puso a Bruto como una furia; se marchó de malas formas, hacia donde tenían atados los caballos. Marco Valerio le gritó mientras se alejaba: 

“¡Maravilloso sitio; mañana estaremos de vuelta! ¡Y disfruta de tu encuentro con la irresistible pero engañosa Venus!” 

“¡Te veré donde siempre, maravilloso imbécil!” 

Me pregunto cuánto duró el ataque de risa a Marco Valerio, que se levantó siguiendo los pasos de su amigo, que se había marchado hecho una divina furia. ¿Podrías decirme que esta amistad, con éstas sus costumbres y actitudes, era la más adecuada para el heredero de un césar? 

Ambos escondían los caballos en la casa de un veterano, un borracho donde los hubiera, que obviamente se entendía sobremanera con Lucio Junio; así evitaban que se los robaran en semejante barrio de mala muerte. Al llegar Marco Valerio al lugar, su amigo ya se había marchado con su caballo. Dejó un sestercio en la mesa del veterano, al que todos llamaban Luperco, nunca he sabido ni querré saber por qué motivos. Por lo que creo, aún no había bebido tanto como de costumbre, y pudo incluso bromear sobre el buen humor de Lucio Junio al pasarse anteriormente por ahí. Marco Valerio le dijo: 

“Está contento porque ha encontrado una mujer que le quiere; pero le va a durar poco el contento. Lo que pasa es que ella es muy mala, y él imbécil.” 

Luperco se río, y ofreció un trago a Marco Valerio, que rechazó amablemente. Se encaminó entonces hacia los establos, donde le esperaba atado su caballo, un hermosísimo purasangre blanco, y montó para desaparecer del lugar. Siempre volvían por un callejón desierto, que conocían perfectamente, y les conducía directamente al final del suburbio. 

Recordarás que, cuando se decidió su traslado a Grecia, se asignó la tutela del pequeño Marco Valerio al entonces pretor, más tarde gobernador de la región: el desgraciado Publio Antonio Saturnino. Que yo consideré la decisión acertada, al igual que lo hicisteis tú y casi todos los que participaron en ella, no hará falta que lo diga, porque era un hombre brillante, sereno y de reputación intachable, además de flexible y comprensivo en general. Lo fue como educador y también, quizá en demasía, como gobernador. El pueblo es traicionero, como lo es por defecto la raza humana, pero además acuciado por la necesidad. Nunca es deseable, pero sí en ocasiones necesario, infundir temor para hacerse respetar por las masas. Yo siempre he considerado en este sentido preferible, aunque no ideal en modo alguno, el modo de proceder de su hermano Cayo Antonio, al que tu mismo predecesor como cónsul, si mal no recuerdo, destituyó como gobernador de Mauritania, por motivo de las matanzas, saqueos y varios abusos que se produjeron durante su mandato. Me reitero en que estas maneras no son en ningún caso las propias del buen gobierno romano; ahora bien, siempre las preferiré a la pasividad y la piedad mal entendida, si bien el anterior cónsul no opinó lo mismo, como ya he dicho.

La hija de este Cayo Antonio, Julia Antonina, también hace tiempo que fue enviada a Grecia, y se instaló, como es natural, bajo la tutela de su tío. Fue así que Marco Valerio y Julia crecieron y fueron educados juntos, y desarrollaron con el tiempo una considerable afinidad. Si bien era hija de su padre, como todo ser viviente que se precie, si a alguien recordaban sus maneras compasivas, agradables y siempre inteligentes era a su tío Publio, de modo que yo siempre la consideré una perfecta compañía para Marco Valerio; incluso en los últimos años, pensé que ella era el necesario antídoto contra Lucio Junio y sus salvajadas. Dicho sea de paso que si no se prohibió que estos dos se siguieran viendo, como yo sin duda habría defendido, es porque nunca lo quiso el procónsul Publio Antonio. 

Aquellas nonas de marzo, supe yo más tarde que al volver del Pireo hacia su casa Marco Valerio se encontró con Julia. Ella era consciente de las travesuras que él se traía con Lucio Junio, y fue a recibirlo a escondidas en las caballerizas. Cuando llegó, él ya había atado el caballo y recuperado su vestimenta romana, y entonces lo abrazó y le dijo: 

“Siempre igual, Marco, hueles a bebida y tienes un aspecto horrible. No quiero saber dónde te ha llevado hoy el chalado de Lucio Junio, y sabes cuánto me disgusta lo que hacéis allí.” 

Marco Valerio le respondió con afecto, acariciando su brazo: 

“Mi querida Julia, preocupada por nuestras gamberradas\dots Tu temor es sin embargo excesivo; si huelo a bebida es por el antro de Luperco, donde no hay otra cosa, y mis malas apariencias se deben a la fatiga de cabalgar. Lucio Junio se ha marchado corriendo, y he tenido que volver sin compañía, lo cual ya sabes que me resulta tedioso.” 

“¡Lucio Junio marchándose solo! ¿Se puede saber qué tramaba ese diablo, que le resultara tan entretenido como para dejar de emborracharse con la chusma por ello?” 

Marco Valerio le habló entonces de la aventura de Lucio Junio con aquella équite desvergonzada y de poca monta, que creo recordar que se llamaba Albina Prima. Obviamente Julia no se sorprendió, ni se indignó contra ella ni contra Lucio Junio, conociendo los desmanes que éste solía traerse, ni tampoco se preocupó excesivamente, sabiendo lo poco de lo que en general le había servido preocuparse por él. Los dos salieron de los establos en silencio, hasta que dijo Marco Valerio:

“Sinceramente, no creo que debamos hacer nada; Lucio Junio obra siempre con tan poco sentido común, y la situación tiene tan pocas salidas, que no llegará a ningún punto que deba preocuparnos. Si la chiquilla no tiene dinero no va a casarse, pero el único dinero es el de los Junios, que obviamente no lo darán para que Lucio se case con ella. Lo único que me preocuparía de él es que en un momento se volviera medio loco y acabara matando a alguien a porrazos.” 

Julia río esto que Marco Valerio había dicho no se sabe si en broma o en serio. Después, en la cena con el gobernador, no se habló nada, que Julia recordara, con el menor interés de cara al caso que nos ocupa. 

\chapter{Idus de marzo}

Las poco sesudas travesuras de Lucio Junio progresaron sin mayores sorpresas; durante unos días él y la équite se veían cada cierto tiempo, ella pedía, él prometía, y entretanto ocurrían cosas que mejor no reproduciré aquí. En un principio no tenía planes, luego se hizo consciente de que necesitaba tener alguno, y entonces decidió que debía casarse con ella, no sé sabe con qué dinero ni convenciendo de qué forma a sus padres. Todo lo que iba pensando, que no era mucho ni muy coherente, lo compartía de inmediato con ella. Siendo Lucio Junio tan fanfarrón, y también lo supieron Marco Valerio y, como consecuencia de ello, Julia Antonina. 

Ella ya no se sorprendía de lo que oía contar, aunque su preocupación por que los otros dos siguieran yendo de aquella manera al Pireo lógicamente aumentó. Las juergas de Lucio Junio eran cada día más alocadas, y acudía cada vez más gente y de peor condición; ni siquiera se molestaba en disimular su origen adinerado, ya que a menudo, estando como una cuba, ofrecía un trago a toda la taberna, incluso pagando en denarios de plata. Se convirtió por tanto en una especie de héroe de esos barrios. Marco Valerio acudía por lo general como testigo, a veces escondiéndose en la esquina de las tabernas para que no lo relacionaran con Lucio Junio en exceso. Es llamativo como incluso él tenía habilidad para relacionarse con aquellos indeseables, a los que hacía disfrutar en conversación, además de que era capaz de hablar perfecto griego sin acento. Lo más probable es que él mismo hubiera previsto el final bochornoso que prometía la situación, y hubiera definido su propio plan de fuga.

Fue entonces cuando llegó el desenlace prometido; ¿cuándo si no en los idus de marzo, malditos para el destino de Roma? En la tasca del día había juegos de dardos y apuestas del circo, por las que Lucio Junio solía perder la cabeza; habiendo perdido veinte sestercios, Marco Valerio decidió alejarlo de la tabla de juegos, lo cual no le costó mucho, ya que bebido como estaba no era propenso a poner muchas pegas. Lucio Junio sin embargo conservaba la consciencia, y luego fue capaz de narrarme con precisión, y sorprendente sinceridad, muchos de los sucesos de aquella noche. Le dijo a Marco Valerio, que lo había llevado a una mesa, en la que se sentó abatido: 

“¿Y si tienes razón, y he perdido la cabeza por esta chica? Los días pasan, solo pienso en ella cuando puedo verla, y en no pensar en ella cuando no\dots ¿Tú crees que me he vuelto loco? Cátulo, tú sabes muchas cosas\dots tienes que saber si me estoy volviendo loco\dots” 

Lucio Junio apoyaba la cabeza sobre sus manos: 

“No creo que te estés volviendo loco —le dijo Marco Valerio—; sencillamente es tarde y deberíamos volver. Además, creo que deberíamos dejar de venir por unos días.” 

“¡Bah\dots!” 

Lucio Junio casi se desmaya después de gruñir de este modo, y entonces, cuando Marco Valerio estaba ya para levantarlo y guiarlo fuera del lugar, se abrió la puerta de la taberna. Afuera reinaba una tarde lluviosa y entraron gotas de agua por la puerta entreabierta; apareció entonces un hombre joven, alto y musculado, malamente vestido, con la cabeza casi rapada y la cara sucia. Se acercó con cuatro amigos al dueño de la taberna, repartiéndoles vinos y olivas. Entonces Lucio Junio, que lo había observado al entrar, pareció despertar de su letargo, y empezó a mirarlo fijamente y con desdén. Entonces le dijo a Marco Valerio al oído: 

“Éste es el payaso de Nepos, un liberto de poca monta; el tipo está como loco por Albina, y se dice que ella también lo quiere, aunque a mí me ella ha dicho mil veces que no es cierto\dots” 

Después de esto perdió el resuello, en el fondo estaba claramente deprimido. Marco Valerio miraba con preocupación evidente hacia todos lados, hasta que lo interrumpió Lucio Junio exaltado: 

“¡Aún hay tiempo para divertirse; pide dos buenas cervezas!” 

“Es el peor momento para ponerse a pedir cervezas; vámonos ya y me hablarás de este Nepos en otro sitio.” 

“¡Bah\dots! ¡Enseguida te hablaré de la rata de Nepos, ahora tú pide dos cervezas!” 

Los clientes alrededor de ellos giraban la cabeza para interesarse por los gritos de Lucio Junio. Incluso Nepos, que ya servido con sus amigos, ocupaba una mesa en otra esquina del local, parecía haber reconocido a su teórico enemigo, y reírse de él con sus amigos. Marco Valerio prefirió mantenerse en la taberna que escapar corriendo, dejando solo a Lucio Junio, al que desde luego ni los truenos de Zeus habrían hecho salir de aquel lugar, ni la magia de Apolo entrar en razón. Fue donde el dueño y le pidió una sola cerveza, dejando un par de dracmas en su mesa, y se dispuso a regresar a la suya con Lucio Junio. Él no perdía de vista al liberto, rozando lo descarado, incluso lo miró desafiante en repetidas ocasiones. A Marco Valerio no le dio tiempo de llegar a su mesa antes de que Nepos se levantara diciendo a toda la taberna: 

“¿Es que ya nadie respeta al bueno de Nepos; habéis olvidado mis servicios y mi incondicional amistad? He oído que un romano loco ha empezado a pagar los tragos de todos los hijos de vecino de este barrio\dots ¡aquí estoy yo, disfrutando el vino que nos da nuestra querida Grecia, y las olivas de Atenas! ¡¿Por qué a mí no me las has pagado, Lucio Junio Bruto; por qué a mí no me da nada la liberalidad de la nueva Roma para con los griegos?!” 

Marco Valerio llegó con la cerveza justo cuando Nepos, habiéndose acercado a ella mientras hablaba, llegaba a la mesa de Lucio Junio. Este último estaba furioso, pero mantuvo su compostura desafiante, mientras le dijo a su amigo: 

“Cátulo, yo te he pedido dos cervezas.” 

Presumo aterrorizado a Marco Valerio, que sólo acertó a decir: 

“Pero yo no quería ninguna\dots” 

“No me has entendido\dots ¡ay con el bueno de Cátulo!” 

Entonces tomó la jarra que le habían traído y la vació con violencia en el rostro del liberto Nepos; éste llevó las manos a los ojos, que se habían llenado de cerveza, lo cual aprovechó Lucio Junio para lanzarle la jarra a la cara, que se rompió del impacto. El liberto estaba aturdido y ciego, no sabiendo por dónde iba; obviamente sus amigos se levantaron enfervorecidos, dispuestos a vengarse del agresor. Fue entonces que a Lucio Junio no se le ocurrió otra cosa que subirse a su mesa, que vació pegando patadas a todos los potajes que había en ella, y sacar la afilada daga que llevaba siempre escondida para decir a gritos: 

“¡Venid a mí y sabréis lo que da la liberalidad de Roma! ¡Venid sucios perros!” 

Los amigos de Nepos, que seguía en el suelo, se lanzaron contra Lucio Junio, que sin embargo no necesitó hacer uso de su cuchillo, ya que los bebedores de las mesas contiguas contentos de la juerga que llevaban teniendo con él (y a su costa) todos esos días, se levantaron en su ayuda. El resto de la taberna, que había llegado demasiado tarde como para disfrutar de los desmanes de Lucio Junio, pareció identificarse con el bando del liberto, interpretando todo aquello como una especie de enfrentamiento greco-romano. Fue entonces que se organizó una verdadera batalla campal. Lucio Junio animaba y dirigía a los suyos desde el puesto de mando en que había convertido aquella sucia mesa de taberna, protegido por cuatro corpulentos borrachos que se habían erigido en una penosa guardia pretoriana. Marco Valerio se escondió en la bodega con el dueño de la taberna, prometiéndole que le compensaría por los destrozos, y siguiendo con el oído los gritos de Lucio Junio, que, mal que pese, eran la única forma que él tenía de seguir los acontecimientos. Descubriendo una salida oculta, salió como una exhalación del lugar.

Así fue que siguieron a mamporros hasta que por fin Nepos volvió en sí. Habiendo sido liberado de las luchas de gladiadores, no le costó tornar la balanza de la pelea en favor de su bando, en cuanto hubo entendido por qué de repente se había alcanzado aquel punto de violencia. Los gritos de Lucio Junio cesaron, porque empezó a centrarse en Nepos, que a su vez lo miraba desde su posición más baja. Se le acercaba entre los golpes y los objetos que volaban a diestro y siniestro por la tasca; Lucio Junio lo miraba desde la mesa desafiante y le apuntaba con el cuchillo, aunque él mismo me dijo más tarde que estaba muerto de miedo ante aquella figura imponente. Nepos se lanzó contra los cuatro borrachos, que poca resistencia ofrecieron y, al mismo tiempo, dos hombres se subieron a la mesa a traición, golpeando a Lucio Junio por la espalda, le arrancaron el cuchillo de las manos y lo apresaron. En aquel instante toda la taberna dejó de pelear, centrándose en lo que sucedía entre Nepos y Lucio Junio. El liberto se hizo con la daga, acercándose a quien la había traído, que estaba sentado sobre un taburete, inmovilizado por un hombre realmente corpulento, a la vez que mugriento, y que no hablaba sino con gruñidos. Con la daga apuntando al rostro de Lucio Junio, dijo: 

“¿Te divierte jugar con dagas, Lucio Junio Bruto? Si es así\dots ¡vamos a jugar!” 

Nepos miró triunfante a su alrededor, recibiendo las glorias de los que habían estado de su parte, cuando Lucio Junio, que no podía moverse un palmo, le gritó: 

“¡Algún día yo te enseñaré lo que es jugar, rata estúpida y torpe! ¡Lo único que te da fuerza es lo que abundan otras ratas como tú! ¡Algún día verás lo que son las ratas; sabréis todos cómo se purgan las ratas\dots!” 

Nepos hizo un gesto al grandullón, que acalló las amenazas de Lucio Junio con una típica llave, seca y dolorosa. Éste resoplaba y se retorcía como podía, cuando Nepos continuó: 

“Algún día igual sabes pelear en vez de hablar más de la cuenta\dots ¿Por qué será que no me sorprende? No es raro que los fanfarrones no sean buenos en el arte del amor.” 

Lucio Junio emitió un gruñido sordo por tener la boca tapada por el gruñón, y entonces sucedió que, justo cuando el repugnante Nepos se le acercaba con la daga, volvió a abrirse la puerta del local. Se había alcanzado una verdadera noche de tormenta, que inundó la entrada de la taberna. Vueltos todos hacia ella, se reconoció la figura de Luperco, seguido de Marco Valerio. Si la figura de Nepos había impresionado a Lucio Junio, no quiero pensar cuánto le impresionó a Nepos la de Luperco. Era de complexión fortísima, y curtido no en la burda arena de los gladiadores, sino en los terribles bosques de Dacia y Germania. Por primera vez desde hacía años, no había bebido nada hasta aquellas horas de la noche, y entró en la taberna en absoluto silencio, con el ceño fruncido y una mirada de perros hacia Nepos. Llevaba consigo una piedra, grande y puntiaguda; la lanzó con una fuerza inmensa contra el hombre que sujetaba a Lucio Junio, acertándole en plena frente, haciéndole caer inerte y ensangrentado al suelo. Luperco se le acercaba, cuando otro de los amigos de Nepos quiso lanzarse contra el veterano. Con suma facilidad, lo mandó al suelo y le quebró la nariz de una patada. Entonces, sin decir palabra, y sin que ningún otro desde luego se atreviera a hacerlo, tomó por la cabeza el cuerpo inerte del hombre, golpeándosela con fuerza dos veces más contra una mesa. Al incorporarse dijo sin levantar la voz: 

“Está muerto.” 

Siguió el mismo silencio sepulcral interrumpido por un trueno; Luperco, manteniendo su humor de perros y su silencio, aparto el cadáver y se fijó entonces en Nepos. Sin decir palabra, le hizo un gesto con la cabeza, señalando la salida, que seguía abierta y mostrando la tremenda tormenta del exterior. Nepos y los otros tres amigos que lo habían acompañado salieron huyendo como una exhalación del lugar. Luperco no se inmutó, y miró entonces de reojo a Lucio Junio, que había quedado liberado, con gesto derrotado. Volvió a señalar la puerta con la cabeza, como antes había hecho. Entonces salieron igualmente todos los que habían peleado a favor de Nepos, y finalmente se vació el lugar, no encontrando los que quedaban muchas ganas de seguir la juerga ni celebrar su triunfo; sólo quedaron el tabernero, el cadáver, el gordo inconsciente, Luperco, Marco Valerio y Lucio Junio. A Marco Valerio no le dio tiempo a dejar doscientos cincuenta denarios de oro en la mesa del tabernero antes de que Lucio Junio se desmayara. 

Lo llevaron él y Luperco al palacio del gobernador, donde Marco Valerio le encontró un aposento cerca del de Julia Antonina. Ella salió a recibirle como de costumbre, pero se encontró con Lucio Junio sin conocimiento, herido y desmarañado, llevado de mala manera por Marco Valerio y un hombre con un aspecto terrible y peligroso, la cual era obviamente la impresión causada a cualquiera que viera a Luperco. Dejaron al demacrado en su alcoba, y Marco Valerio despidió al veterano lleno de agradecimiento, y dejándole una bolsa tintineante, todo ello para el mayor asombro de Julia Antonina. Ella miró entonces a Marco Valerio de modo inquisitivo, él le apartó la vista y se acercó a Lucio Junio. Veía cómo entreabría los ojos de vez en cuando. Los dos se apostaron en la cama, cada uno a un lado, y Marco Valerio dijo: 

“¿Por qué haces esto, Lucio?” 

Lucio Junio sólo emitió un leve e incomprensible sonido. Entonces Julia le dijo suave pero claramente: 

“¿Qué has hecho; ni siquiera a mí vas a decírmelo?” 

Él, sin incorporarse, abrió los ojos, que se le llenaron de lágrimas, y dijo casi sin voz: 

“Ella\dots creo que ella no me quiere\dots ella quiere a la rata\dots y a mí me quiere para reírse de mí\dots hoy lo he sabido; me he dado cuenta de que está cansada y ya no me quiere\dots” 

Y después volvió a cerrar sus ojos llorosos; Julia Antonina y Marco Valerio se miraron llenos de estupor. Lo dejaron durmiendo y se marcharon a sus aposentos. 

A la mañana siguiente, ambos se despertaron realmente temprano, para que Lucio Junio no fuera descubierto. Lo llevaron a su casa, donde Julia dijo que había ido a montar a caballo cerca del Muro Largo y que había sufrido una caída, por lo que decidieron que se quedara a dormir con ellos. Obviamente los padres de Lucio Junio no eran idiotas, y además conocían el carácter de su hijo, pero ¿quién iba a dejar de creer a la elocuente e intachable Julia Antonina, excelente sobrina del procónsul? Además, el aspecto demacrado de Lucio Junio bien pasaba por el de alguien caído de un caballo, incluso atropellado por un carro. Lo dejaron allí al cuidado de sus sirvientes hasta que despertó; curiosamente, después se inventó la misma mentira que Julia Antonina, fuera por telepatía o por haberla entreoído en su convalecencia. 

Ella y Marco Valerio estaban obviamente llenos de preocupación. Éste último se veía claramente superado por la situación, según ella misma me dijo más tarde; andaba en silencio, cabizbajo y se apoyaba con frecuencia la cabeza entre las manos. Me dijo también que Marco Valerio era ciertamente empático y de naturaleza sensible, lo cual no me sorprendió mucho, pero además me dijo que en ocasiones se mostraba incapaz ante los cuadros graves como aquél, lento a la hora de proponer soluciones y torpe al llevarlas a cabo. Digamos que, superado por aquel orden de las cosas, lo único que hacía era sufrir. Finalmente deliberaron sobre el asunto, pero no decidieron nada. Evitar las juergas de Lucio Junio, si no se daba él mismo cuenta de la necesidad de ello, iba a ser tarea imposible; el tipo era testarudo a más no poder, doy fe de ello. 

Sin embargo, el día siguiente, Marco Valerio debió de verlo con bien distinto semblante. Julia Antonina se alegró sobremanera simplemente al cruzarse con él por la mañana; al parecer, Marco Valerio era de naturaleza realmente peculiar, y si bien por su lentitud cualquier problema le hacía cavilar en exceso, digamos que el doble que a una persona ordinaria, después su solución también era el doble de profunda y bien pensada. Es así que a Julia le alegró verlo de buena gana. En el almuerzo se acercó a él en silencio, y sin necesidad de preguntarle por el asunto comenzó él mismo: 

“¡Julia Antonina, qué gusto verte por aquí; toma unas olivas, que nos han recogido esta misma semana, y son excelentes!” 

Ella me dijo que efectivamente lo eran, pero que no le dejó de disgustar que Marco Valerio se pusiera a apreciar la excelencia de las olivas justo entonces, lo cual no dejaba de corresponder a un cierto carácter altivo y socarrón. Entonces le ofreció también una jarra de vino: 

“¡Cuánta suerte tienes de haber venido\dots el mejor vino de Grecia! Y el garum\dots ¿dónde está el garum?” 

Julia lanzó una mirada furiosa a Marco Valerio, quien obviamente tuvo que captarla, por si pretendía hacer sobre el vino de Calcídica la misma disertación que sobre las olivas de Atenas; él sin embargo mantuvo su sonrisa burlona, en cualquier caso no sin llegar a donde pretendía Julia Antonina: 

“El vino es espantoso\dots porque nos hace olvidar nuestros problemas, y eso nos aleja de poder arreglarlos\dots ¿lo entiendes, Julia?” 

“¡¿Qué demonios pretendes que entienda\dots— Julia Antonina verdaderamente hecha una furia— me estás tomando el pelo de veras?!... ¡Ayer Lucio Junio casi se mata\dots y pudo haber sido una masacre!... ¡Y hoy seguirá medio borracho tumbado en sus aposentos\dots y a ver cuándo demonios se levanta! ¡¿Se puede saber dónde pretendes llegar hablando tanto de vino\dots has estado toda la mañana bebiendo vino?!” 

Marco Valerio se río notoriamente, mirando a Julia de un modo ciertamente extraño en él. Su tono era sarcástico y burlón, lo cual no era, me dijo, raro en él, pero sí en una situación drástica como aquella; sin embargo, como en aquel momento no estaba para preocuparse por Marco Valerio, sino por Lucio Junio, no le dio mayor importancia. Él efectivamente tenía algo pensado, lo cual evitó que Julia Antonina le organizara la merecida escena: 

“En general es preferible que la gente no se atiborre a vino —Marco Valerio terminaba de reírse—\dots porque a los borrachos se les limpian los sesos, y se acaban quedando sin ideas\dots el vino pudre la mente y acaba con lo mejor de uno; sin embargo\dots ¡qué no daría yo a cambio de que a Lucio Junio se le borraran las ideas que tiene en la cabeza!” 

“¿Y bien?\dots” 

“Pues que\dots ¡tenemos que enseñarle lo que es la paz romana! Lucio Junio pasará una semana yendo de la farra, a los juegos, a las luchas\dots dentro de cinco o seis días se despertará, ni él sabrá en qué cuchitril de mala muerte, medio borracho y alelado del todo, y entonces se le habrá olvidado lo de esa estúpida équite, y ese repugnante liberto, y él mismo se avergonzará de haberla deseado de tan mala manera\dots ¡hoy lo llevaremos al circo, Julia Antonina!” 

Ésta quedó obviamente extrañadísima, si no lo estaba ya del todo, a raíz de esta última intervención de Marco Valerio; ahora bien, el momento requería preocuparse por Lucio Junio, que bastantes problemas había causado, y de que la solución de Marco Valerio parecía ajustarse a lo que requerían las circunstancias, por extraño que pueda parecer. Julia pensó que quizá sólo estaba un poco alterado de cuánto había sucedido aquellos días, pero en principio aún en sus cabales, y así es como Julia Antonina ignoró por completo su comportamiento, y accedió a la solución propuesta\dots ¡ante sus ojos estaba el delirio, el preludio de la desgracia de la locura! Ella pudo darse cuenta de que Marco Valerio no estaba en sus cabales, y comprender que la cuestión de Lucio Junio lo había superado por completo; habría podido decirle que ella era capaz de resolverlo del todo, que él podría olvidarse\dots ¡podría haberlo hecho, cambiando la suerte de Roma, pero no lo hizo! 

En este singular estado en el que se encontraba Marco Valerio, salieron ambos a la calle, dispuestos a seguir lo planeado; él iba decidido y a paso firme, y ella lo seguía insegura y maquinalmente: pasarían primero por casa de los Junios, después llevarían a Lucio al circo y al final a pegarse la mayor farra que recordara, para acabar dejándolo medio muerto en cualquier esquina. Esta última condición era la más segura tratándose de Lucio Junio, y lógicamente estaba ausente de peligro si se realizaba en una villa de buena familia, en vez de en fonduchos perdidos en barrios de mala muerte. Nada podía preverse, salvo que al despertarse estaría tan perdido que se habría olvidado de Albina y de cualquier otra tontería al respecto. 

Sé sin embargo que Julia Antonina no estaba satisfecha con la idea. En primer lugar no era muy partidaria de las orgías, ni pensaba que fuera bueno ser vistos en una. Después había que contar con una reacción favorable de Lucio Junio, que bien era capaz de montar una escandalera borracho y deprimido. Por último, no podía dejar de extrañarse del comportamiento de Marco Valerio. Así todo, no tuvo otro remedio, falta de ideas y de réplicas, que atenerse por absurdo que fuera a lo que éste había pensado.  

Mayor se hicieron su asombro y sus miedos, en todo caso, al encontrarse con Lucio Junio. Llegados al recibidor de la villa de los Junios, empezaron a oír gritos desde el interior: 

“¡Ajá—reconocieron la voz de un Lucio Junio perjudicado—\dots! ¡Os lo dije\dots han venido! Y ellos solos\dots ¡sin que yo les dijera nada! ¡Ajá\dots ja!” 

“¡Cátulo! —apareció otra voz, ésta desconocida. ¿Y ésa no es\dots?” 

“¡Es Julia Antonina, la sobrina del procónsul\dots ¿también ella viene?! ¿Dímelo Bruto, ella también viene? ¡¿Eh, Julia, tú también vienes?!” 

Aparecieron entonces los tres en el vestíbulo. El segundo en hablar había sido Gayo Junio Terco, primo de Lucio Junio, que al parecer había bebido menos vino que su primo. El tercero había sido el auténtico sinvergüenza de Quinto Tulio Básico, una deshonra en todos los aspectos para su buena familia. Todos abrazaron llenos de alegría a Marco Valerio, que parecía divertirse, y a Julia, que cada vez observaba más llena de horror. Si la idea de Marco Valerio de atiborrarlo de vino parecía mala y absurda, peor aún era el que exactamente lo mismo se le hubiera ocurrido a Lucio Junio, y más aún le aterrorizaba que Marco Valerio empezara a tener las mismas ideas que su colega descerebrado. Entonces entró en la sala el hijo del prefecto, Claudio, hombre joven pero rechoncho y desfavorecido, además de limitado, desde luego apuntando tan malas maneras como su padre; estaba borracho como una cuba. Llegó, casi se cae al suelo y gritó, sin mirar a nadie: 

“¡¡Mentecatos!!” 

Entonces Claudio se retorció por el suelo de la risa que se daba a sí mismo. Divertido con semejante panorama absurdo, Lucio Junio dijo al resto del grupo: 

“¡Decidido queda\dots se acabaron todas las tonterías! ¡Primero al circo\dots y luego iremos a casa de esta bola de sebo\dots lo de hoy en la Villa Claudia dicen que va a ser para acordarse\dots hip!” 

“¡¡Mentecatos\dots ah\dots hip!!” 

Entonces marcharon todos, Lucio Junio adelante intentando sostener y guiar a Claudio, después Marco Valerio con Gayo Junio, y detrás Quinto Tulio no paraba de hacer guiños acercándose a Julia, que lo contemplaba con horror y andaba llena de nervios, intentando adelantarse junto al resto del grupo: 

“Así que ha sido todo una broma —dijo sarcástico Gayo Junio a su primo—, lo de esa caballera, Albina Prima, o se llamaba\dots” 

“¡Ah\dots! ¡No\dots! ¡Una broma no; fue una estupidez! Una estupidez, una simple y llana estupidez\dots ¡pero ya me he arrepentido! ¿Entiendes\dots? Fue una tontería, un capricho tonto, pero ya me he arrepentido\dots ¡¡del todo; y ahora que no se hable más del asunto!!” 

“¡Ja —Gayo le contestó entre risas, unidas a las de Quinto Tulio y a un sonido incomprensible que emitió Claudio—, ¡no decías eso hace dos días\dots!” 

“¡Sí\dots sí\dots —continuó Tulio Básico— entonces presumía el bueno de Bruto de su finísima conquista\dots ajá!” 

“¡¡Valientes mentecatos\dots!!” 

“¡¡Sois odiosos, imbéciles todos!!\dots ¡¿Cómo habré ido a acabar con semejante panda de imbéciles\dots?!” 

“¡¡Mentecatos!!” 

Lucio Junio se calló lleno de furia después de la enésima interrupción de Claudio, pero pronto volvió a tomar la palabra: 

“¡Ah\dots! ¡¿Qué más me da en el fondo si sois imbéciles?! Hoy sí que vamos a divertirnos, y entonces me dará igual si soy un estúpido, o si todos vosotros sois imbéciles\dots ¡bah\dots qué me importa a mí nada de eso!” 

Julia observó la escena llena de estupor, aunque al menos le sirvió para quitarse de encima al insufrible Tulio Básico, y así observar medianamente bien a Lucio Junio y a Marco Valerio. El primero parecía que obviamente fanfarroneaba, aunque al menos era consciente de que todo aquel asunto de la caballera había sido una tontería, y se esforzaba en olvidarlo a gritos. Marco Valerio siguió la escena en silencio, pero dando muestras de estar divirtiéndose de lo lindo con aquellos tipos tan excéntricos. A Julia no le extrañó esto en demasía; no era raro que se divirtiera, que se relacionara con tipos de cualquier clase, tampoco en el fondo que riera sus gracias, que se mostrara burlón y sarcástico, como en la comida anterior, pero había algo, fuera en su gesto, en su mirada, algo que ella no podía describir, que le resultaba extraño, sumamente extraño\dots 

La multitud se agolpaba alrededor del circo, en la lucha por hacerse con un sitio, pero Claudio avanzaba al frente del grupo, apartando a la gente a empujones, de forma que tal circunstancia poco les obstó para llegar a ver los juegos a la hora. Tenían asientos separados, por grupos de tres. Lucio Junio ofreció el palco familiar, pero Julia, no queriendo ser vista con aquellos borrachos, le miró con tal cara de perros que por suerte se le quitó la idea. Se sentaron Marco Valerio, Julia y Gayo Junio abajo, y los otros tres, que estaban incluso en peor estado que al salir de la casa de los Junios, unas cuatro filas por encima. Julia corroboró durante los juegos, sin mucha sorpresa, que Gayo Junio era el más inteligente de los tres amigos de Lucio, y además casi no estaba borracho. No queriendo estar en silencio, y además deseosa de escuchar hablar a Marco Valerio, cuyo estado le intrigaba sobremanera entonces, dijo: 

“Entiendo entonces que también vosotros estabais al tanto de lo que sucedía entre Lucio Junio y aquella caballera.” 

“¿Al tanto? —Gayo Junio respondió entre risas—. ¡Lo tiene que saber ya toda Roma!” 

“No quiero saber cómo lo están pasando los Junios\dots ni mi pariente Julia Severa\dots!” 

En efecto, Julia era pariente lejana de la madre de Lucio Junio. Mientras hablaban comenzaron los juegos: 

“¡Ah\dots ellos sólo saben una parte de la verdad\dots creen que ha sido una amante sin más! Ni se han imaginado que el tipo estaba incluso dispuesto a casarse con ella\dots ¡habrían muerto del susto de saberlo!” 

“¿En serio quiso llegar tan lejos?” 

“A mí nunca me dijo tal cosa —dijo Marco Valerio, evitando quedarse en silencio—, en todo caso.” 

“Desde luego; un día me confesó sus intenciones: decía que iba a robar siete talentos de oro de la cámara de la familia, porque era todo lo que podía esconderse entre la ropa sin ser visto. Después yo mismo tenía que inscribirlos como casados en los archivos de la ciudad, porque a menudo trabajo en el ayuntamiento, y me habría sido fácil hacerlo. Entonces su plan era escaparse lejos, no sé dónde\dots pero lejos; esto me lo confesó a regañadientes después de mucho insistir. En un primer momento yo me negué, pero entonces me dijo que se iría aun sin estar casado\dots desde luego había perdido el norte\dots” 

“¿Y es posible —prosiguió Julia— que haya olvidado todos estos proyectos de la noche a la mañana?” 

“Viéndolo ayer —respondió Marco Valerio—, es imposible que los haya olvidado\dots” 

“Lo bueno es —cerró Junio— que al menos se ha decidido a olvidarlos.” 

Se oyó precisamente a Lucio Junio y a Claudio, cuando pasó el líder de la carrera de cuadrigas: 

“¡Borracho\dots tú no tienes que ganar! ¡Borracho tramposo! ¡¿Dónde están los azules: los AZULES?!” 

“¡Mentecatos!” 

No hace falta ni decir que al final perdieron los azules, y que Lucio Junio y Claudio perdieron una fortuna, lo cual, por otro lado, les importaba bien poco. Pasado el jolgorio de la carrera, Julia, que aún no había satisfecho sus dudas, se dirigió a Marco Valerio: 

“¿Marco\dots tú eres, sin embargo, el que más cerca ha estado de Lucio Junio\dots? ¿Tú crees posible que se haya quitado el asunto de la cabeza, y con qué motivo?” 

“Bueno\dots ¡ayer casi nos matan; quizá ello le haya hecho pensar que ha llegado demasiado lejos! Si yo hubiera aparecido un minuto más tarde con Luperco, estos juegos serían en honor del ilustre Lucio Junio Bruto\dots ¡muerto de mala manera en una sucia tasca!” 

Julia Antonina estiró el cuerpo, como alejándose de Marco Valerio, que se encontraba a su izquierda, y se llevó las manos a la cara como señal de espanto. Gayo Junio, a la derecha, leyendo el gesto de Julia, quiso consolarla, como diciendo que, por suerte, ya se había acabado el asunto. Julia respondió casi sollozando: 

“¿Pero cómo? ¡¿Cómo se ha llegado a este punto?! Una vida ilustre, llena de felicidad\dots ¡¿cómo decides estropearla de esta manera\dots!? Es suicida\dots ¡suicida!” 

“Quizá —espetó Marco Valerio—, pero también tiene su punto de natural.” 

Gayo Junio y Julia miraron a Marco Valerio al mismo tiempo con un asombro absoluto: 

“¿Cómo que de natural?” 

“Ah Julia\dots ¡te gusta demasiado el orden romano; no puedes entender lo que otras personas buscan! Yo he pasado las tres últimas semanas con Lucio Junio de juerga en juerga, y he aprendido a entenderlo, en cierta manera\dots desde luego no era una cuestión de lujuria, pero diría que ni siquiera de amor; quería vivir de otra manera, sufrir de otra manera\dots” 

“¿Sufrir de otra manera\dots —preguntó Gayo Junio— qué demonios significa eso de sufrir de otra manera? ¿Estás diciendo acaso que tenía razón mi primo?\dots ¡yo tampoco entiendo nada!” 

“Ah Terco; no entendéis nada\dots y no; no le doy la razón; sólo intento comprenderle, aunque he de confesar que me ha llegado a dar cierta envidia —esta palabra envidia hizo que Gayo Junio y Julia se miraran llenos de extrañeza—. A Julia le gusta el orden romano; a ti quizá no te interese tanto el orden romano\dots pero es igual; partamos de él. Concedo que la paz romana efectivamente es feliz, pero esta felicidad no tiene nada de interesante. Haz lo que te dicen tus padres, y te buscarán un buen matrimonio, sonríe a los ancianos y tendrás un puesto alto, divierte e imita a los jóvenes y tendrás multitud de amigos\dots sé un poco más inteligente que la gente que te rodea y triunfarás\dots ¡y si no rodéate de gente menos inteligente! ¡La felicidad está a la vuelta de la esquina!\dots pero los hombres no están hechos para conformarse con la felicidad, y por eso habrá cada día más gente estúpida como Lucio Junio, gente a la que no le gusta la paz de Roma. Lo que pretendía Lucio en el fondo, estoy seguro, era vivir una aventura peligrosa\dots ¡eso es, una aventura peligrosa!” 

“Mi querido Marco —Julia no daba de sí de su asombro—\dots ¡pero qué estás diciendo! Todo esto es irresponsable, caótico, inmoral\dots es hasta blasfemo\dots ¡pero ante todo me parece absurdo! ¡La paz romana feliz, la paz romana mala! ¡¿Pero desde cuándo es algo malo la felicidad?!” 

“Ah\dots ¿no son contrapuestos la felicidad y el dolor; no aparece la una gracias al contraste con el otro? Sin dolor no hay felicidad, pero somos unos cobardes y llamamos felicidad a la ausencia de dolor\dots ¡que en el fondo lo es también de felicidad!” 

Julia estaba realmente asustada, y no sabía qué pensar: 

“Pero\dots ¿qué es eso de la aventura peligrosa? ¿Por qué no se alistó si quería aventuras peligrosas? ¿Por qué no se hizo sangrar las manos si quería sufrimiento\dots? ¡Todo esto es un absurdo!” 

“¡Ya lo creo! —Junio también quiso tomar parte en tan extravagante conversación—. ¡Anda con Marco Cátulo; vaya una forma de hablar\dots vaya una tontería de forma de hablar! ¿Las barbaridades de Lucio Junio se nos van a presentar como aceptables por ponerles otro nombre: de aventura peligrosa?\dots ¡eso no quiere decir nada!” 

“Ah\dots Julia\dots Terco\dots no entendéis lo que digo—ni siquiera se enfadaba, respondió con absoluta naturalidad—. Pelear con bárbaros no es peligroso\dots en el sentido que digo; el verdadero mal es el dolor, y el verdadero peligro arriesgarse a él, pero mucha gente no puede asumir casi ningún sufrimiento; no al menos en la época nuestra\dots bajo el reinado de la paz romana. Ni tampoco el sufrimiento del que hablo es el de los enfermos ni los doloridos ni los caídos en combate; hablo del dolor de la derrota\dots ¡sí, eso es, del dolor de la derrota! Yo creo que Lucio Junio ha sido tonto pero no ciego, un día vio en profundidad, quizá se imaginó toda su vida rodeado de formalidades y estupideces\dots y decidió apuntar más alto; entonces tuvo una ambición distinta, y por enfrentarse a ella no le importó asumir el dolor de la derrota, ¡y por eso al menos, debemos entenderlo; incluso podríamos llegar a admirarlo!” 

“¡Admirar al ambicioso de mi primo\dots! ¡¿Pero cómo puede ser alta la ambición de manosearse con una équite limitada y desvergonzada?!” 

“Bueno; quizá su ambición era sólo\dots la de tener una aventura peligrosa.” 

Gayo Junio calló poco convencido, pero riendo con la última intervención, que debió de encontrar ingeniosa, de Marco Valerio; el estado de éste siguió preocupando, incluso más que anteriormente, a Julia, que sin embargo no quiso seguir hablando del asunto, y dijo: 

“La verdad es que tiene mérito, Marco, que hayas conseguido entenderle. Es lógico, de todos modos, si habéis pasado tanto tiempo juntos\dots pero estemos contentos de que, incluso a él, lo de ayer le resultara un exceso insostenible en el tiempo, y se esfuerce en dejar todo esto de lado. ¡Creo que lo peor ha pasado, y debemos estar contentos!” 

Julia dijo esto último mirando de reojo a Marco Valerio, y explorando nuevas reacciones esclarecedoras de su estado, pero él simplemente asintió y sonrío. Ella se centró por tanto en los juegos, los cuales por otro lado detestaba profundamente, y consideraba una pérdida de tiempo. Arriba, Lucio Junio, Quinto Tulio y Claudio estaban verdaderamente montando un espectáculo, ridículo incluso para los lamentables comportamientos que suelen darse en las carreras. Uno insultaba a los corredores, lanzaba objetos, el otro escupía a las demás gradas, se oía ese absurdo y continuo grito de “mentecatos”, y naturalmente cada vez que pasaba el vendedor de bebidas se le compraba toda la mercancía. Por suerte transcurrieron los juegos, gracias a la paciencia del resto del circo, sin mayores incidentes, lo cual es llamativo, porque en situaciones menos graves sé que ha sido llamada una patrulla, y yo mismo he llegado a tomar represalias importantes contra este tipo de gamberros insufribles. Perjudican la paz romana, y también la paz y el orden de algunas mentes notables\dots ¡como habían hecho con Marco Valerio, no lo dudo! 

Transcurrieron así los juegos, y los tres de la fila de abajo, gracias a la previsión de Julia, se marcharon algo antes del final, antes de la desagradable acumulación de gente en la salida. Los otros tres salieron sin embargo con el montón, lo cual no les costó gracias a los empujones de Claudio, que organizaron un verdadero tumulto de indignación, del que sin embargo, ágiles entre la multitud, lograron zafarse. 

Llegaron entonces a Villa Claudia, donde, como preveía Claudio, no iba a haber más que dos o tres esclavos, ya que el prefecto se había llevado a los demás consigo a no sé qué otra orgía que había organizado, ¡creo que en la misma prefectura! Pido de paso que ordenes el destierro de semejante sabandija, cuya deposición, como es natural, ya he ordenado en cuanto supe de sus prácticas deleznables. Así es que aparecieron en la vacía villa, Lucio Junio sentó a todos en el salón—recibidor, que era enorme, llamó a unos músicos que había contratado para la ocasión, pidió a Claudio que organizara a los esclavos que habían quedado, para que sirvieran bebidas y aceitunas. Con el paso de las horas, fueron entrando más amigos a los que Lucio Junio había llamado, u otros que ni conocía nadie, pero se habían enterado del ruido y querían igualmente divertirse de lo lindo. Lucio Junio puso a otros a jugar a los dados, que parecían de poca monta, porque a ellos ya no los conocía, diciéndoles: “en esta casa se juega a los dados, y con apuestas altas, y si no no se viene a esta casa para nada\dots ¿habéis entendido, borrachos?”. Los otros obviamente no tuvieron otro remedio que ponerse a jugar a los dados. Entonces sacó pieles de leones y cocodrilos, traídos del país del Nilo, y cascos de legionario que se guardaban en la casa y obligó a otra gente a ponérselos, y a estos les puso a jugar a que los legionarios perseguían a los cocodrilos, los cocodrilos a los leones y los leones a los legionarios. Obviamente entre borrachos no hace falta dar órdenes cuando se tiene una simple idea. Así, quedaron Marco Valerio, Julia, Quinto Tulio y Gayo Junio en los asientos del centro de la sala, alrededor de los cuales jugaba a perseguirse toda esta agrupación de borrachos absurdos; en un extremo de la sala estaban los jugadores de dados, que la habían empezado a pegarse, en otro los músicos, en el otro los esclavos aterrorizados y en el último, el más apagado claramente, se había quedado dormido Claudio encima de un charco de sus vómitos\dots y por todas partes estaba Lucio Junio, disfrutando de semejante colección de absurdos. 

Julia Antonina miraba en todas direcciones, no por curiosidad sino preocupada de que alguien conocido la viera participando en semejante cuadro, y si seguía ahí era por preocuparse de Marco Valerio. Éste sin embargo no hacía sino conversar amablemente con Gayo Junio, conversación a la que se unieron Julia y Básico. 

Pasaron varias horas así; habían decidido disfrutar de la noche como espectadores, que es como ellos me lo describieron más tarde, riendo las gracias absurdas de Lucio Junio y el espectáculo que estaba organizando, evitando meterse en los previsibles líos, y entretanto hablando de estupideces triviales, en opinión de Julia y de Gayo Junio. La música iba decayendo, la mayoría de los jugadores de dados se habían ido arruinados, indignados o apalizados, y los disfrazados se iban escondiendo, por parejas o en grupos más grandes, en diferentes rincones de la casa, puedes imaginarte con qué intenciones; ¡para que veas la calaña que se juntó en esa casa! Lucio Junio llegó a decirme, no sin sorna y un cierto orgullo juvenil, que no sabía si había sido más escandaloso lo que organizó aquel día o lo que llevaba haciendo las semanas anteriores. 

Así transcurría la noche hasta que Julia se sintió verdaderamente sin fuerzas, y se levantó del sofá donde había estado conversando con Marco Valerio y Gayo Junio, hacia otro que entonces estaba libre. Entonces estaba agotada y pensaba dormirse allí mismo. Ellos dos se mantuvieron en su sitio, hablando de poesía y bebiendo unos tragos. Sin embargo, Básico se levantó al igual que Julia, y la siguió hasta el sofá donde se había colocado. Ella se había sentado y él se puso a su frente; seguía terriblemente borracho: 

“Sabes que a mí —le dijo susurrándole al oído—\dots a mí no me dicen nada las chicas serias y respetables\dots” 

Él esperaba quizá que Julia respondiera algo, pero cansada como estaba no dijo nada, además de que, desde luego con buenos motivos, a él lo despreciaba profundamente: 

“Pero hoy\dots hoy me he dicho\dots hoy me he dicho que\dots ¡vaya con la buena de Julia Antonina!” 

Entonces el muy desvergonzado se tiró de mala manera encima de la pobre Julia, quien desde luego bastante había aguantado, durante todo aquel día rodeada de borregadas y sandeces; ella emitió un grito ahogado, quizá lo que le hacía falta para despertarse, se zafó de él de un tortazo y gritó: 

“Ah\dots tú\dots ¡¡desgraciado!!” 

Se sopetón se detuvo la música, y Marco Valerio y Gayo Junio se levantaron como un resorte hacia Quinto Tulio, al que prendieron cada uno de un brazo; incluso Claudio se despertó de su lecho perfumado, se acercó al lugar tambaleante, y gritó: 

“¡¡Mentecato!!” 

Julia miraba a Básico con cara de perros, y éste decidió desaparecer de la villa sin despedirse de nadie. Lo mismo hizo la gente que había traído Lucio Junio, y que aún estaba en el salón, considerando quizá que se había hecho tarde, también se fueron los músicos y los esclavos, que salieron corriendo\dots Lucio Junio espetaba a todos aquellos que se marchaban, como diciéndoles “pero\dots estaba siendo una maravilla\dots como os marcháis de esta villa maravillosa\dots”, y así hasta que todos se marcharon. Entonces se sentó abatido en el centro del salón. En aquel preciso momento, se levantó Julia como un resorte y dijo: 

“Esto ha sido deleznable\dots ¡deleznable! ¡No quiero permanecer en esta casa de locos, borrachos y sinvergüenzas ni un minuto más\dots ni uno! ¡¡Me voy de aquí ahora mismo!!” 

Marco Valerio y Gayo Junio se levantaron con ella, como haciendo amago de irse también, pero entonces intervino Lucio Junio, desde luego indignado del todo con que tuvieran que acabarse sus diversiones: 

“Ah\dots ¡¡ah\dots!! ¡Me abandonáis todos\dots! ¡¡Todos!! ¡Ha quedado claro que aquí no puede uno divertirse de ninguna manera\dots así que cabalgaré\dots ya lo creo que cabalgaré: hacia las buenas tascas del Pireo, donde los jornaleros y los esclavos fugados todavía no han empezado a beber\dots!” 

Todos lo miraron entonces llenos de preocupación; él se marchó corriendo por la entrada de la villa y desapareció de su vista. Esta escena, unida a la de Básico y al cansancio que la amilanaba, hicieron que Julia olvidara casi del todo las rarezas de Marco Valerio durante el día; se acercó a él y le suplicó que siguiera a Lucio Junio, que evitara que cometiera alguna otra estupidez. Se instaló allí, en un aposento de la misma Villa Claudia, y Marco Valerio salió diligente en busca de su amigo. Gayo Junio, debido al estado de Claudio, se ocupó entre tanto de la organización de la casa. 

Marco Valerio se hizo con el mejor caballo de la villa, y marchó al galope en busca de Lucio. Éste debía de haber alcanzado un ritmo altísimo, porque incluso saliendo del centro de la ciudad no hubo manera de alcanzarlo. Reinaba la noche sobre el campo que conectaba la vieja Atenas con el Pireo, y Marco Valerio se guiaba por las teas en lo alto de cada torre del Muro Largo, además de la suya propia. Fue mediada la muralla que se encontró con otra que avanzaba igualmente hacia el puerto; entonces le gritó con todas sus fuerzas. Efectivamente, era Lucio Junio, que se detuvo y llevó su caballo al encuentro de Marco Valerio: 

“Ah\dots ¡Cátulo; yo sabía que tú no ibas a fallarme\dots! ¡¡Siempre nos quedará Cátulo\dots!! Ahora vamos\dots ¡la noche no ha acabado\dots nos esperan, ya lo creo que muchos nos esperan\dots!” 

Justo cuando Lucio Junio arrancaba, Marco Valerio le adelantó con su caballo, poniéndose justo a su frente, y el caballo de Lucio tuvo que hacer una verdadera pirueta para esquivarlo. Lucio Junio miró extrañado a su amigo, y luego ordenó a su caballo que continuara, pero éste, agotado por el esfuerzo anterior, amedrentado quizá por la oscuridad de la noche y no habiendo esperado el peligroso encuentro con Marco Valerio, desobedeció a su jinete, y se quedó parado: 

“Hoy —aprovechó el momento de pausa para hablar Marco Valerio— no\dots hoy no despertaré a Luperco\dots ¡no lo despertaré! No sé a quién pretendes ofender hoy, ni que tasca vas a poner patas arriba\dots ¡pero yo no podré ayudarte! He venido hasta aquí simplemente para decirte esto.” 

“Ah\dots ¡sea así Cátulo!—dijo disponiéndose a espolear el caballo de nuevo—.” 

“¡Espera!” 

Lucio Junio se volvió de nuevo hacia Marco Valerio, que adoptó un cierto tono sarcástico, como anteriormente en la comida con Julia: 

“¿Y no te parece injusto que siempre decida el bueno de Bruto? Ah\dots ¡como Lucio Junio Bruto está en crisis, pues hay que ir a las tascas que dice Lucio Junio Bruto, pelearse con los borregos que miran mal a Lucio Junio Bruto, llamar a borrachos curtidos en el arte de la lucha para que saquen a Lucio Junio Bruto de un apuro\dots ¡eh\dots es que nadie tiene para mí una crisis amorosa como la de Lucio Junio Bruto!” 

“Ah\dots ¡Cátulo; eres insoportable del todo! ¿Adónde quiere ir el señorito caprichoso?” 

Marco Valerio sonrío, espoleó a su caballo, ya descansado como el de Lucio Junio, y le dijo a éste que lo siguiera. Fueron ambos a casa de Luperco, donde desmontaron, se cambiaron de ropas y tomaron sus clásicas dagas, necesaria precaución, que llevaron escondidas entre la ropa. La noche empezaba a dejar paso a las primeras luces de alba, de modo que uno en cierto modo podía distinguir por dónde andaba. Marco Valerio avanzaba al frente, con paso decidido y demostrando conocimiento de las callejuelas. Por fin, doblada una esquina, Lucio Junio ya supo hacia dónde se dirigían\dots y al de poco apareció ante ellos el Mediterráneo: 

“Creo es la paz del mar lo que ahora te hace falta, incluso a ti, Bruto.” 

Lucio Junio asintió, y fueron recorriendo el paraje, hasta el mirador desde el que el anterior día habían estado mirando el muelle del repugnante Avidio Régulo, que entonces estaba vacío, contrastando con el ajetreo del anterior día. Al llegar allí, Lucio Junio notó cómo empezaba a hacerle efecto la última copa de vino que había tomado al salir de Villa Claudia. Justo entonces se aproximó alguien al muelle; la luz del amanecer se hizo con el lugar, haciéndole perfectamente visible: era una muchacha, vestida humildemente con ropas griegas. Lucio Junio, al verla, dijo medio en broma, ya nuevamente bajo el efecto embriagador del dios Baco: 

“¡Mírame a mí, el amante despechado\dots yo llorando como un idiota\dots como si no estuviera el mundo lleno de mujeres espectaculares!” 

Dijo más tarde que esto lo había dicho borracho, pero que no recordaba que la chiquilla fuera tan especialmente hermosa. Debía de tener el pelo oscuro, contrastando con una piel curiosamente pálida, lo cual es raro en un esclavo, incluso después del invierno; entonces se había sentado en el muelle, mirando con cierta gracia el mar. También es cierto que Lucio Junio es un chico vanidoso hasta la médula, y es raro que reconozca de buenas a primeras que de una muchacha que es hermosa, a no ser que crea haberla conquistado. Incluso creo recordar que, aparte de en esta ocasión, por extraño que pueda parecer, nunca más tuvo ocasión de verla. 

Marco Valerio, que parecía algo ausente y ni tan siquiera había reparado en ella al salir, dijo de repente: 

“Me ha mirado.” 

“Eh\dots ¡¿qué?!” 

“He dicho que me ha mirado y que es preciosa.” 

Lucio Junio no tuvo tiempo de saber lo que se le estaba diciendo; se desmayó justo entonces por el efecto de la bebida; cerró sus ojos viendo, fija en él, la mirada misteriosa, sarcástica y burlona de Marco Valerio \dots

\chapter{Calendas de abril}

Sucedió que al día siguiente Lucio Junio se despertó, dolorido y aturdido, en el palacio del procónsul, recordando fragmentos inconexos de lo sucedido la noche anterior. Sabía que había estado con Marco Valerio en el Pireo, de lo que dedujo que éste había tenido que transportarlo no sé sabe cómo hasta su residencia, pero nada recordaba aún del extraño encuentro que habían tenido con aquella esclava griega. Pensó en Albina Prima con repugnancia, lo cual lo llenó de contento de sí mismo; tal pensamiento y tal contento lo condujeron a levantarse, precisamente en busca de Marco Valerio. Deambuló por el palacio hasta por fin encontrarlo en el jardín trasero. Estaba leyendo un volumen, que Lucio Junio no reconoció, mientras picaba pan y aceite de oliva. Parecía tener un humor estupendo, en contraste con el aturdimiento y las bruscas maneras que traía Lucio Junio; hacia éste se volvió, al darse cuenta de que había llegado. Entonces ambos se miraron un instante, Lucio Junio pasmado y aturdido, y el otro sonriente y burlón, hasta que Marco Valerio no pudo aguantarse y entró en un ataque de risa, que interrumpió Lucio Junio: 

“¡Eh\dots tú! ¿Se puede saber de qué te ríes ahora?” 

“De nada\dots ¿durmió bien el señor?” 

“¡A mí qué me importa si he dormido bien! Tengo que hablar contigo, de cosas muy importantes, ¿entiendes?\dots ¡muy importantes! ¿Y por qué estás tan contento? ¡Te juro que no entiendo por qué estás tan contento\dots!” 

Marco Valerio enrolló sus pergaminos y miró a Lucio Junio: 

“No lo recuerdas; interesante cuanto menos\dots ¿puedo preguntar cuántas cosas recuerdas, de lo que pasó por la noche? Obviamente no muchas\dots” 

“Recuerdo que debemos disculparnos con Julia por bastantes cosas\dots ¡pero no recuerdo qué; te juro que quiero acordarme, pero no me acuerdo!” 

“Julia, nos gritará, se enfadará y dirá que somos unos sinvergüenzas. Llevo toda la mañana pensando qué decirle\dots pero, en fin, ya se nos ocurrirá algo. ¿Cuál era, si se puede saber, el asunto tan importante del que ibas a hablarme?” 

“¡Ah, por supuesto! Verás: mi enfermedad de amor alcanzó límites insospechados; ¡no te imaginas cuántas noches en vela ni cuántas lágrimas le he dedicado a esta chica, cómo mi tristeza se esfumaba cuando ella aparecía, para volver después al instante, de qué modo estaba presente siempre en mi cabeza, pero todo eso fue en vano, ¡y ya está olvidado! Mi enfermedad de amor está curada, ¡ya lo creo que lo está! Pero no es de amor que aparecen todas las enfermedades: también he sido engañado y ofendido; y todo eso bien distinta cura requiere; ¿lo adivinas, verdad?\dots ¡desde luego que tienes que adivinarlo!\dots¡¡Se requiere, por el honor de mi familia y de mi patria, una venganza!!” 

“Así que ahora te preocupas del honor de tu familia y de tu patria, después de tratado de ponerlas en ridículo de todas las formas que has ido encontrando\dots” 

“¡Ah\dots Cátulo; dime que me entiendes!” 

“Desde luego; en su momento lo mejor fue enamorarse de la mujer equivocada y de forma estúpida, poner patas arriba Grecia, y Roma, o lo que demonios hubieras querido pasar por encima, y ahora hay que olvidarlo todo; y tanto lo has olvidado que necesitas vengarte\dots” 

“Ah... ¡valiente imbécil!” 

No siguieron con la conversación. Lucio Junio desayunó y se tumbó exhausto. Después de un rato, se marchó a su casa. Para despedirse, ambos se dirigieron a la salida: 

“Hay una cosa que me gustaría saber—dijo Marco Valerio—: ¿para tu venganza particular habrá que derramar sangre?” 

Lucio Junio se quedó un momento mudo, pero luego resopló diciendo: 

“Pues aún no lo he pensado\dots pero seguramente sí.” 

“Sea.” 

Lucio Junio siguió el camino un rato en silencio, pero entonces se paró pensativo y dijo: 

“La verdad, a mí me gustabas porque usabas la cabeza\dots ¡pero ahora que eres también un irresponsable, también me haces gracia!” 

Marco Valerio no respondió, mirándolo con su habitual sonrisa maliciosa. Llegaron a la salida y dijo: 

“Ve a Villa Claudia; encontrarás a Julia, y entonces descríbele de todas las maneras que puedas lo lamentable que fue lo de anoche. Dile que no volverás a verte con Tulio Básico, que estás indignado, y que yo también estoy igualmente indignado. Pero obviamente ni se te ocurra decirle que al alba fuimos al puerto, ¿me entiendes?” 

“Estupendamente; ¿algo más desea el César Augusto?” 

“Nada, estaré ocupado\dots a la espera de cuando empecemos con tu venganza, claro.” 

Se despidieron, marchando Lucio Junio hacia Villa Claudia, como había dicho. Creo que desde aquel momento estuvieron más de una semana sin verse. Marco Valerio se entretuvo con unos de filosofía, que al parecer le tenían bastante ocupado. Lucio Junio, por su parte, tampoco estuvo muy visible; se encerraba para elaborar bocetos de lo que pretendía que fuera su plan de venganza: un día enviaría no sé qué matones por ahí, después prendería fuego a la casa de Albina, después secuestraría a Nepos y lo sometería a crueles torturas, que él había aprendido de sus contactos en Oriente\dots Obviamente semejante ejercicio de creatividad requería cierto descanso, y Lucio Junio se lo tomaba con esporádicas juergas con los pocos que quedaban en Atenas tan brutos como él, a excepción de Tulio Básico, ya que había prometido a Julia, en sus propias palabras de disculpa, “no volver a mezclarse con semejante indecente”. Tampoco hace falta decir que las juergas de “descanso” ocupaban la amplia mayoría de la jornada de Lucio Junio, que dedicaba a sus maquinaciones a lo sumo diez o veinte minutos al día. 

Pasada esa semana, Marco Valerio volvió a encontrarse varios días con Lucio Junio, simplemente para hacer sus habituales gamberradas. No sucedió nada resaltable hasta el primero o segundo de abril, cuando Lucio Junio por fin avanzó su proyecto de venganza. Se paseaban por el templo de Erecteo de la Acrópolis, habiendo dicho Marco Valerio que aquel día se las bastarían tomando el aire y sin recurrir a la violencia ni a las borracheras de Lucio Junio. Fue entonces que éste, sin que hubieran hablado mínimamente del asunto, dijo: 

“No derramaremos la sangre de ella, pero sí la de él. Ella debe sufrir, viendo como él muere.” 

“¿Se puede saber qué dices?” 

“¡Ah —Lucio Junio se tiró encima de Marco Valerio—\dots! ¡Estoy diciendo que él morirá: que mataremos a la rata de Nepos\dots con mis propias manos! ¡Y ella tendrá que verlo!” 

“Ajá\dots —Marco Valerio no pudo contener una sonrisa llena de ironía—. ¿y cómo pretendes hacerlo, Sócrates del Ática, Apolo del Olimpo?” 

“¡Ah —Lucio Junio gesticulaba como si le fuera la vida en ello—\dots! ¡Pues aunque no lo creas, tengo algo pensado; lo he organizado todo para que él muera\dots con ella delante! ¡La sangre de la rata dará de beber a la arpía\dots! ¡Recibirá lívida un baño de roja sangre!” 

“¿No has pensado que ella igual se ríe de él como hasta ahora se había reído de ti? Igual le traería sin cuidado, o incluso le haría gracia, que abrieras las tripas de Nepos delante de ella. ¿Y si encima después te prenden los guardias, o te linchan los vecinos de su barrio\dots? ¡La matarás de risa; igual en eso consiste tu venganza! ¿Me estás diciendo en serio que no habías reparado en este aspecto hasta ahora?” 

“¡Maldito seas Valerio Cátulo —Lucio Junio no podía contener sus gritos—\dots tienes que decir algo con respecto de todo!” 

Entonces empezó a jurar, dando vueltas por todas partes, lanzando piedras y gritando a cualquier lado, hasta el punto de que algunos niños que habían ido allí a jugar a los pañuelos se marcharon aterrorizados. Todo ello fue, naturalmente, para el mayor regocijo de Marco Valerio. 

“¡Ajá\dots! ¡Todo cuadra a la perfección con mi plan, y no me vas a decir esta vez que no! ¡¿A que es cierto que, si te encontraras con cualquiera de los involucrados, sean Albina, Nepos, o cualquiera de ellos, no te reconocerían?!” 

“Totalmente; sólo me he encontrado con esa gente el día que, tan sabio que eres, casi te descuartizan en aquella tasca, y yo juraría que ninguno de ellos reparó en mí. De hecho, estaba hablando con aquellos sembradores medio borrachos y luego\dots en fin, juraría que no me reconocerían.” 

“¡Ahí está el asunto, y además hablas perfecto griego, tanto que se suelen pensar que eres griego en ti mismo! ¡Serás mi infiltrado! Necesito que te introduzcas en el círculo de máxima confianza de Albina Prima, ¡y después me dirás exactamente cómo y de quién tengo que derramar sangre delante de ella! ¡¡Y entre medias nos cargaremos a Nepos, eso dalo por sentado!! ¡Así se haga y que no se hable más del asunto!” 

Lucio Junio estaba al borde de un ataque de nervios; lanzó su mirada febril y penetrante a Marco Valerio. Éste le correspondió con un ataque de risa: 

“¡Anda con el loco sanguinario! ¡Y yo resulta que soy ahora un espía romano en territorio extranjero!” 

\textbf{Cohesionar}

Fue a unos tres o cuatro pasos de Lucio Junio, en el momento en que éste se iba incorporando, que se paró, tomó aire, y volvió a hablar ya con más tranquilidad: 

“Así que —decía mientras se tornaba hacia Lucio Junio—\dots así que te gusta hacer el imbécil\dots ¡y quieres que yo haga lo mismo\dots!” 


“¡Ajá\dots nublado por el ciego y tirano Cupido! ¡Así es como quiero que hagas el imbécil!” 

“¡Tú querrías que hiciera el idiota movido por la ceguera y los sentimientos más estúpidos!\dots pero no será el caso; lo haré todo sencillamente porque quiero\dots ¡me apetece y ya está! Pero no discutamos sobre los motivos de mis actos; creo que tenemos un maravilloso proyecto del que ocuparnos. ¿Quieres que me cuele en la casa de Albina disfrazado de ánfora? ¿Podría ser rellena de aceitunas?” 

“¡Ah\dots eres estúpido, no te aguanto! Esta vez tengo sin embargo un plan que será infalible\dots  ¡y ni siquiera tú podrás burlarte de él! Escúchame bien —dijo hablando deprisa, como si quisiera demostrar la rapidez y la amplitud de su ingenio—: a casa de Albina Prima va cada dos o tres días un pobre barbero, emigrado de Tiro. El tipo se dedica a arreglarle el pelo de mala manera; hasta un tonto podía hacer eso, pero ella le paga porque le gusta pensar que tiene servidumbre.” 

“¿Quieres que lo matemos?” 

“¡Ah\dots pero qué estás diciendo! ¡Siempre intentas tomarme el pelo! El tipo debe de no estar contento con su vida, lo cual no me extraña. Entonces lo sobornaremos para que se cambie de barrio, ése es el primer punto. Y le diremos que no se moleste en pasarse por casa de Albina Prima a despedirse. Será entonces que tú te presentarás ahí, hablando griego con el acento que quieras, y dirás que te manda a sustituirle. ¡Entonces te las arreglarás para aprenderlo todo sobre ella, ganarte su absoluta confianza, y así completaremos nuestro plan de venganza!... ¡rodéate de tus enemigos y acabarás venciendo, tal y como dijo Julio César!” 

No hace falta que diga que no sé en absoluto de dónde demonios había sacado Lucio Junio que Julio César hubiera pronunciado palabras semejantes. Conviene que te aclare en este punto que, a pesar de haberse enamorado estúpidamente de ella, Lucio Junio nunca fue más que un amante de Albina Prima, y de hecho casi no sabía nada de ella ni de su más profunda intimidad, ni tampoco de sus verdaderas amistades. Con quién se relacionaba lo supo de ella por encima o de lo que le dijeron sus otros contactos en Atenas. En su brutalidad habitual, a Lucio Junio le habría bastado con liarse a mamporros con Nepos, y probablemente haber acabado por salir perdiendo. Hacía días que, sin embargo, le dominaban los delirios de lo que él entendía por una venganza impresionante y sofisticada, todo lo cual era aún más estúpido si tenemos en cuenta lo poco que entendía de sofisticación. Eso explicaba tanta necesidad de lo que él llamaba conocimiento de su víctima. En el fondo, todo residía en su obsesión por Albina Prima. Su monomanía en busca de su amor se había canalizado en la necesidad de odiarla y vengarse de sus afrentas, pero en realidad Lucio Junio no quería vengarse de ella, simplemente necesitaba tenerla presente de continuo. 

Después Lucio Junio contó a Marco Valerio lo que sabía sobre la familia de la susodicha; su padre había muerto hacía años, y su madre se quedó loca por completo, hasta el punto de que acabó mudándose a Dalmacia con su otro hijo, un niño de tres o cuatro años. Albina se quedó entonces sola con la casa familiar en Atenas, y con bastante renta para vivir holgadamente, entre la herencia de su padre y los esporádicos envíos de su madre. Su reputación entre los bien nacidos ciudadanos de Roma era tan mala que no había una familia decente en la que no se prohibiera el trato con ella, pero esto no se debía a su situación, ni tan siquiera a su rango, sino a su comportamiento bajo y lascivo, del cual Lucio Junio no había sido el único, aunque quizá sí el más estúpido, afectado. Decía todo esto rápido, como sin darle importancia, contrastando con la profunda atención de Marco Valerio: 

“Y, en fin, qué voy a decirte\dots tampoco te estoy pidiendo gran cosa\dots mañana te presentas ahí, te imaginas que eres uno de esos histriones de teatro, o del estilo, y analizas el panorama de Albina, la sobrina revoltosa, la criada ladrona y salchichera, y la bellísima esclava: esa sufrida joya en semejante pedrusco de vileza\dots ¡sólo tienes que combinar análisis e histrionismo, cosas que por cierto se te dan muy bien!” —concluyó. 

Se callaron entonces, Marco Valerio pensativo, y Lucio Junio deseoso de oír una respuesta, por un instante. Aquél, después de moverse de un lado a otro, y con el rostro muy serio, como si estuviera enfadado, miró a su amigo: 

“¡¿Se puede saber de dónde te has sacado todo esto?!” 

“Ajá\dots ¿y eso a ti qué te importa?”  

“Bueno, bueno\dots sí\dots tienes razón; lo importante es lo que uno sabe, y no cómo ha llegado a saberlo. ¡Cúmplase entonces tu voluntad, decídase nuestra suerte\dots! Francamente, es un todo un asunto muy divertido; sí, ése es el motivo de que me guste\dots ¡es divertido!” 

Dicho esto, casi se cae al suelo de la risa, para el estupor de Lucio Junio: 

“¡¿Será posible que te hayas convertido definitivamente en un chalado?! ¿Pero qué tontería dices de que es divertido, y por qué se supone que te hace tanta gracia?” 

Marco Valerio se le acercó entonces, sonriéndole y mirándolo fijamente, con aspecto efectivamente de estar chiflado: 

“¡Ah\dots! ¡No entiendes nada, Lucio Junio\dots! No estoy hablando de lo que es divertido, como se uno divierte tirando piedras al río, o jugando esos estúpidos juegos\dots ¡hablo de una diversión última, eterna, primigenia! Verás, Lucio Junio —dijo entonces en voz baja, como si de un gran secreto se tratara lo que estaba diciendo—, hace tiempo que le he estado dando vueltas\dots no a un plan, como has hecho tú, ni a una idea, o a un objeto, que es como normalmente piensa la gente, sino a todo: a la cuestión, aquello sobre lo que hay que pensar, ¿entiendes, Lucio Junio? Porque no se puede hacer geometría, si no se sabe lo que es una línea recta, del mismo modo que no se puede vivir, ¡si no se entiende lo que es la vida! Es entonces que he llegado al fondo de la cuestión\dots ¡seré arrogante por hablar de esta manera! He visto todo lo que hay que ver\dots ¡que es lo mismo que lo poco que hay que ver! No hay más que vacío, en el fondo de todas las cosas, así que\dots ¡así que no debo tener reparos en ponerlo todo patas arriba! Divertido\dots peligroso\dots es un juego, porque no sé dónde están los juegos, si no hay diversión ni hay peligro\dots ¿vamos a divertirnos con un juego, y me dices que no es éste un juego lo suficientemente divertido?” 

Lucio Junio miró a Marco Valerio con una expresión que combinaba pena y espanto, dando a entender que se había dado cuenta del punto hasta el que se había vuelto loco. Quizá evaluando esta expresión, Marco Valerio se apartó, y sonrío, dando a entender que le había gastado a su amigo una especie de broma mal interpretada. Hubo un rato de tenso silencio, Marco Valerio sonreía y Lucio Junio miraba extrañado. Debes notar, ilustrísima, en todo caso, un elemento esencial sobre el carácter de Lucio Junio: si bien puede tener buenas ideas, o alcanzar percepciones sensatas, es imposible que ninguna de ellas se le mantenga en su cabeza de tres libras durante más de medio minuto. Si por ella pasó la perfecta visión de la locura de Marco Valerio, que ya empezaba a dar sus muestras evidentes, ésta se difuminó con la misma rapidez que había llegado. Y así las cosas, sin reparar demasiado en el asunto, cambió de expresión y concluyó: 

“Bueno, tú lo has dicho\dots ¡qué importarán las razones!” 

Y así siguieron tan tranquilamente el resto de la tarde, pero ya no se dijeron más locuras, sino estupideces, en detrimento del interés con respecto a lo que trato de narrarte. 

Llegados a este punto, es imposible que Marco Valerio no se hubiera dado cuenta de lo absurdo de la idea de Lucio Junio. Que es inadecuado en un joven patricio romano disfrazarse de peluquero y colarse en una casa teñida de mala fama, no hará falta que yo lo diga, aunque tampoco que a Lucio Junio y Marco Valerio eso les traía sin cuidado. Sin embargo, era improbable que consiguieran mantener hasta el final la discreción, y más aún que Marco Valerio se las bastara para introducirse en el círculo más íntimo de Albina Prima, cosa que ni siquiera Lucio Junio había hecho en su momento. En caso de conseguirlo, a saber qué se le iba a ocurrir hacer entonces a Lucio Junio, ¡y cuál iba a ser el resultado en caso de llevarlo a cabo! Si Lucio Junio deliraba y se dejaba llevar por la monomanía, está claro que Marco Valerio no lo hacía menos. La diferencia entre ambos, siendo la desgracia el resultado en ambos casos, la marcaba la más elemental inteligencia; así como Lucio Junio no había medido nunca de ninguna manera sus actos, y no había previsto las consecuencias de lo que tenía pensado más allá de los pasos más inmediatos, más si cabe con respecto a la forma de comportarse de Marco Valerio, éste seguramente era consciente de hasta los más reservados detalles de la forma de la que iban a desarrollarse los acontecimientos, pero, perturbado como estaba, no los consideraba con la debida angustia, sino que seguramente se regocijaba imaginándoselos. El uno no olía la sangre, y el otro, tocando a las puertas de la locura, se prestaba a saborearla.

%========================
\chapter{Poco después de nonas de abril}

La casa de Albina Prima era lo único salvable de un barrio de mala muerte; se notaba que en su tiempo su padre fue medio-rico. Era una despejada mañana de primavera, y ella esperaba en el vestíbulo, en la parte de abajo, sentada en un diván y con rostro sereno y una mirada maliciosa, que eran lo habitual en ella. Se abrió la puerta y entró Cibaria, una cocinera de Avidio Régulo. 

Esta Cibaria resultaba ser una prima segunda de Nepos, y en casa del mercader se dedicaba básicamente a pisar uvas y pelar pescados. Eso sí, siendo tan infame como su primo, a veces se lleva algunas de esas uvas, pescados y también quesos, especias, y cualquier cosa que fuera robando en la cocina de la casa del mercader. Los que ocultaban lo que ella iba robando, a cambio, claro está, del sustancioso porcentaje que corresponde, eran justamente Albina Prima y su primo segundo Nepos. 

Cuando entró en casa de Albina, su gordura y la cuerda con la que ataba su túnica grisácea y maloliente recordaban de hecho a los fardos de pescado que se cargaban en el puerto; y más aun teniendo en cuenta que bajo las ropas llevaba dos libras de aceitunas, un pescado y una sarta de salchichones. A modo de saludo, puso todo eso sobre una mesa, y después miró a Albina con el gesto del niño que ha cumplido con sus deberes en la escuela: 

“¡Ajá\dots qué hermosura!” —dijo ésta con ironía mirando a la cocinera. 

Entonces se levantó de su sitio y fue a buscar algo, que resultó ser un denario de plata: 

“Aquí tienes —dijo mientras se lo lanzaba con un punto de desdén sin mirarla. La cocinera se arrojó al suelo torpemente a recoger la moneda, que ni siquiera había caído donde debía— ¿Sabes, Cibaria?, he sido demasiadas veces injusta contigo; pienso que te he hecho sentir menos que el resto\dots ¡hoy te quedarás con los demás, te prometo que no te volveré a decir que me haces pasar vergüenza!” 

“¡Con los demás—Cibaria hizo sonar su voz intensa y estridente—, ah\dots qué alegría!” 

“Sí\dots serás lo que le hace falta al grupo; mi peluquero, ese dálmata estúpido, se ha marchado, sin despedirse siquiera, simplemente dejando una nota, ¿te lo puedes creer? —Cibaria puso cara como de pretender sorprenderse—. Pero parece que ha encontrado un amigo tonto que puede sustituirle. Un idiota de Epiro, o un lugar del estilo\dots ¡qué más me da, un tonto u otro! A éste le pagaré menos, eso sí\dots”

Justo entonces apareció Sempronia, la anciana criada de Albina, acompañada de dos jóvenes. Una de ellas estaba ostentosamente vestida, era alta y espigada y sonreía con desdén; la otra iba por detrás, más modesta, pálida, la mirada interrogativa, sencilla y curiosa de unos bellos ojos claros. Albina despidió a Sempronia, que se marchó con los manjares traídos por Cibaria, y comenzaron los saludos. La primera de las recién llegadas era Avidia Emilia Paulina, sobrina del conocido Marco Avidio Régulo. La arrogancia y el desdén que desprendía eran el perfecto reflejo de su más profundo carácter; acostumbrada a la ostentación y al lujo, creció mimada en un entorno gobernado por la laxitud más deleznable, aunaba un exagerado aprecio de sí misma con una irremediable debilidad de carácter. Si frecuentaba la casa de la desvergonzada Albina, donde se cometía adulterio, donde se robaba a su tío sin el menor reparo ni disimulo, era porque, en sus palabras, Albina era la única persona que podía entenderla y estar a su altura, lo cual era sin duda arbitrario y estúpido. Lo que sucede es que la soberbia, cuando no se junta con la inteligencia, siempre acaba por recurrir al autoengaño y la arbitrariedad. 

Emilia se sentó en un diván enfrente de Albina, quedándose de pie quien la acompañaba, que era su dama de compañía, llamada Nereida. ¿Debería acaso denostar a la quizá más relevante causa de nuestra desgracia? Quizá debería hacerlo, pero sería injusto con la realidad. Nereida era una muchacha modesta, sensata, inteligente, y, además, era una esclava, ¿qué razones y qué posibilidad tenía para evitar el cariz que tomaron los acontecimientos? En aquel momento Albina tomó la palabra: 

“¡Mi querida Emilia, qué bien tenerte aquí ---dijo con socarrona altanería---!  Un imbécil de Epiro o de no se sabe dónde me viene recomendado como peluquero\dots ¡a ver de qué palo se trae, yo no voy a meter a cualquier idiota en esta casa! Quiero que lo juzgues como sólo tú sabes, querida\dots” 

Emilia sonrió complacida y en ese justo momento entró Sempronia con los salchichones de Cibaria preparados en un plato: 

“¡Ah, qué maravilla! Emilia, fiel amiga, ¿es necesario que te acompañe siempre esta esclava tuya? No me fio de ella, y podría delatarnos.” 

“¡Ja\dots Nereida! ¿A quién va a delatarnos, si no habla con nadie\dots? ¿Verdad que eres inofensiva?” 

“Bueno, bueno, deberá una confiarse —continuó Albina sin dar tiempo a contestar a la aludida. Se levantó entonces poniéndose enfrente de la esclava, diciéndole desafiante: —Y bien\dots ¿qué opinión te merecen las trastadas de tu ama y de su amiga desgraciada? ¿Nos consideras un escándalo?” 

“Ni te molestes en preguntarle esto\dots” 

“¡Que conteste!” 

Emilia miró por tanto a Nereida como diciéndole que hablara: 

“No me corresponde a mí opinar sobre lo que haga mi ama ni la gente con la que se mezcla —terció esta última.” 

“¡Lo ves, no opina ni se preocupa de nada: un absoluto cero a la izquierda!” 

Albina se sentó de nuevo, riéndose perversa con Emilia. Entró otra vez Sempronia, anunciando la llegada de un joven: 

“¡Que entre! —dijo Albina—. Así que es joven\dots” 

Podrás imaginarte lo que pensé yo cuando me contaron lo que sucedió entonces. Entró Marco Valerio en la estancia, vestido con unas ridículas y viejísimas ropas plebeyas griegas, y que también se había dejado una barba descuidada, el pelo oscuro desmarañado. Claramente estaba interpretando un papel; sus gestos y movimientos eran absurdos y exagerados, en contraste a su talante tan delicado y elegante habitual. Desconociendo este último, me imagino que las presentes tomaron aquel histrionismo como el carácter genuino del hombre que acababa de aparecer. Marco Valerio se quedó enfrente de los dos divanes, mirando alternadamente, con una extraña sonrisa, a Emilia y a Albina Prima: 

“Y\dots ¿a quién debo el honor de la invitación? —dijo abruptamente.” 

Ambas se miraron en silencio, extrañadas, como es natural, de semejante aparición. Albina no quiso permitir que su extrañeza se interpretara como duda o debilidad, así que le espetó a Marco Valerio con autoridad: 

“¡Será posible\dots qué te importa a ti quién te haya invitado; estás aquí para cumplir una función! ¡Mira mis manos resecas, y mi pelo abandonado y descuidado! ¡Con todo el trabajo que tienes y te atreves a preguntar de quién es esta casa; oh, qué osado!” 

Marco Valerio no le estaba haciendo el menor caso; se dedicaba a pasear la mirada por la estancia mientras se le hablaba. Volvió su fija mirada sobre Albina, que tenía una verdadera cara de perros, y le dijo: 

“Es verdad\dots ¡que estáis horribles! Bah\dots ¿cómo queréis que os lo diga? —decía sin aguantarse algunas risotadas— ¡Estáis horribles, dais vergüenza\dots todo esto es patético! ¿Con qué gusto se ha decorado esta estancia? Estos son sin duda muebles de época\dots ¡de la de los dorios; y ya eran feos entonces! Y estos mosaicos, ¡son horrorosos!”

“Pero ¡cómo es posible esto\dots—Emilia se levantó hecha una furia, sintiéndose humillada—esto es una infamia, una verdadera infamia! ¿Cómo permites esto en tu casa? ¡Ahora mismo echa a este individuo!” 

Albina estaba muda y no sabía qué decir; Marco Valerio, por su parte, volvió a no darse por aludido ni lo más mínimo. Pasó la mirada por delante de Emilia, igual que si no hubiera dicho nada, y la dirigió a Albina: 

“Ajá; ¡así que eres tú la señora de la casa!” 

Cibaria y Nereida, que hasta este momento estaban pendientes de otras tareas, no pudieron evitar levantar sus cabezas mirando al visitante. Esta última seguramente sentiría algo pena por la situación, y quizá había encontrado a Marco Valerio simpático en su histrionismo. Cibaria, llena de vileza, deseaba sin duda ver una escena. Ésta, sin embargo no se produjo; Albina volvió de repente en sí, recuperó su semblante risueño y altanero, y dijo a las demás:

``No, no\dots me divierte. ¡Que pase! Dime, ciudadano —se levantó de su diván desgastado—, ¿cómo tengo el gusto de poder llamarte? ¿Cuál es tu origen?''

No me fue fácil entender por qué Albina no echó en ese momento a Marco Valerio, y menos aún, sobre todo ante el riesgo de que pasara eso mismo,  por qué este último decidió aparecer con tales formas en aquella casa. Con el tiempo he entendido, sin embargo, que los aires señoriales de Albina escondían naturalmente su verdadera falta de agallas. Quizá temía profundamente que si hubiera expulsado al recién llegado éste se hubiera negado a marcharse, y que entonces su autoridad hubiera quedado en entredicho. Por otro lado, Marco Valerio era consciente, aunque no con muchos detalles, de que Albina, si bien en el fondo para aprovecharse de él, había desarrollado en su momento una cierta predilección hacia Lucio Junio. Por tanto, quizá aparecer con maneras exageradas y extravagantes como las suyas podía servir, de alguna manera, para acercarse a Albina Prima, más que si hubiera aparecido simplemente para cortarle el pelo y darle las gracias por su confianza. No hace falta que diga que Emilia, con su débil carácter sometido a Albina, asimiló su nuevo criterio, y pareció divertirse con el recién llegado como la que más. Marco Valerio no tardó, por su parte, en responder:

``¡Ah\dots me preguntas por mis oscuros orígenes, y por mi nombre despechado y caído en desgracia! Me llamo Nikandros, como mi infortunado padre, quien fuera el más rico y respetado mercader del puerto de Apolonia, esa joya bañada por el Adriático, allí por donde pasan, en dirección a la gloriosa gran urbe, el trigo de Farsalia, el aceite de Atenas, el hierro de Macedonia \dots ¡incluso el cuero y la carne en salazón, de la tan lejana Dacia! Así crecí yo, en el mayor lujo y refinamiento, aprendí las más intrincadas sofisticaciones, y me eduqué por los mejores maestros de la ciudad. Ni el gusto de los mejores vinos, ni el toque de especias de las tierras más remotas, ni las proporcionadas esculturas de los más reputados artesanos, ni las precisiones de la gramática, ni tampoco los galimatías de la geometría y del estudio de los astros, eran un misterio para mí\dots pero por favor, pongámonos manos a la obra, ya que no quiero aburriros con mis historietas, carentes de la menor trascendencia\dots''

Casi al unísono, Emilia y Albina, que en absoluto esperaban al nuevo peluquero demostrando tal nivel de sofisticación ni tan misterioso pasado, le insistieron en que continuara. Cibaria, decepcionada al no haberse producido una bronca, pasó a otras tareas, mientras que para Nereida el nuevo visitante había cobrado considerable interés. En cualquier caso, no sin hacerse derogar un poco, continuó Marco Valerio con su disparatado relato, ante la insistencia de sus anfitrionas llenas de curiosidad:

``Bueno, habré de compartiros mi pasado, al ver que tanto interés os produce\dots aunque he de confesar  lo doloroso que resulta para mí recordar sucesos tan aciagos ---por lo que me contaron, la capacidad de demostrar tristeza por parte de Marco Valerio contando estas fábulas fue asombrosa---. A los catorce o quince años, no había yo conocido, como os he dicho, más que la comodidad de la opulencia, ni que el calor del hogar paterno. Ya en esos tiempos comenzaba yo a aprender el oficio de mercader, considerando mi familia la necesidad de que continuara los pasos de mi padre. Justo entonces empezamos a conocer sus aspectos más oscuros, bien lejanos a su reputación intachable. Un día, a mi madre le avisaron de haberlo visto en los bajos fondos, en compañías que mejor no menciono. Ella, por supuesto, no creyó al criado que nos lo había comentado, expulsándolo de inmediato de nuestro servicio. Sin embargo los sucesos venideros habrían de sacarla de su incredulidad\dots al de pocos días ---Marco Valerio (o debería decir Nikandros de Apolonia, tan bien que estaba interpretando su papel) dijo esto entrecortado, como si fuera a romper a llorar--- apareció en casa ido, oliendo a vino, y con la cabeza y las manos manchadas de sangre\dots''

Es cuanto menos llamativo que, al menos Albina y Nereida, que no eran en absoluto tontas de remate, no notaran nada extraño en el relato de Marco Valerio. Evidentemente, el barbero que suele aparecer en casa de la mayoría de la gente no resulta ser el hijo de un magnífico mercader venido a menos. Parecía descartable sin embargo la alternativa más plausible, a saber, que el visitante se tratara de un truhán de poca monta tratando de impresionar a sus anfitrionas, viéndolo a todas luces atesorado de un gran entendimiento. Ante esta circunstancia, la realidad (el que se tratara de un patricio romano importantísimo haciendo una trastada altamente peligrosa y totalmente absurda) era sin duda infinitamente más descabellada, y no creo que a ninguna de las presentes se les pasara jamás tal cosa por la cabeza.  





%----------------Descripción de Nereida-------------
La chiquilla debe de ser hija de una esclava tracia, o dacia, que la tuvo con algún griego, que al parecer tenía que ser de pelo muy negro y ojos claros. El tipo se escapó del meollo sin dejar ni rastro, y se sabe poca cosa de él. En cualquier caso, tenía que ser muy listo, porque la chica, aunque no se ocupa nadie de ella, debe de ser muy espabilada, y de su madre no debe de haber salido semejante genética. Ha debido de aprender a leer ella sola, y hablas con ella como si fuera una académica; ¡inusitado, desde luego que eso es inusitado! 

Lo que pasa es que el condenado Avidio Régulo no ha sabido nunca muy bien qué hacer con ella, después de la paliza que le dio a su pobre madre por tenerla. No se sabe si por remordimientos de conciencia, o viendo que la chiquilla es muy capaz, pero no le gusta la idea de tenerla anudando redes o pelando pescado. Y como tiene una sobrina, o sobrino-nieta, que es también muy jovencita, la ha unido al séquito de ésta, que no deja de tener así un elemento algo exótico del que presumir, ¡es increíble, todo increíble! 


\chapter{Idus de abril} 
Unas mazmorras al uso habrían sido letales para el pobre loco ambulante, tan debilitado físicamente tras su marcha hacia la ciudad. Además, la presencia en una prisión de personajes semejantes conducen a escándalos y amotinamientos casi sin excepción. Por tanto dispuse que se le llevara a una villa, abandonada hace siglos, en medio de los olivares más allá del Muro Largo, donde quedó bajo la más estricta vigilancia. También quise que se le diera alimento en abundancia, posibilidad de pasearse y aseo, además de la atención médica que le resultara necesaria. Seis semanas después del lamentable espectáculo que provocara, allí seguía sin causar incidentes de ninguna clase. 

Así las cosas, decidí hacerle una visita a mediados del mes de abril, cuando la agenda de los otros muchos asuntos que tenía pendientes me lo permitió. Por un lado, quería determinar hasta el último detalle quién era ese hombre, su procedencia y la causa de que hubiera acabado en un estado tan deplorable. Por otro lado, y toda Atenas pensaba lo mismo que yo entonces, su aparición se trataba sin duda de un mal augurio, de un desfavorable designio de los dioses, y era en él dónde debía averiguarse la naturaleza del mismo; ¡cuánto más tarde supe, sin embargo, cuán profunda habría de ser la desgracia que auguraba!  

Cuando llegué a la villa, me recibió el capitán de la patrulla a cargo del lugar. Me dijo que el loco ambulante se había empachado, y que tras ello dormía plácidamente la siesta. Casi todos los días se había negado a bañarse, pero aquél era, nadie entendía por qué, una excepción; según me dijeron, para mi fortuna. También me dijeron que parecía haber olvidado quién era; nunca preguntaba ni hablaba con nadie, aparte de lo más elemental de cada momento, y se dejaba llevar maquinalmente y sin rechistar por cada una de las instrucciones de sus vigilantes.  

Habían de prepararse, pues, todos los detalles, sin dejar nada al azar. El soldado más valiente y fiable de la casa acompañó al sujeto a la habitación más recóndita, donde se le hizo entrevistarse conmigo. Nada debía oírse fuera de este espacio, y ninguno, aparte de los tres presentes, debía saber lo que se dijera en esa entrevista. No se encontró mejor lugar que el cobertizo, y tuvieron por desgracia que haber más testigos: una vaca coja, que no habría podido moverse de aquel lugar sin organizar un esperpento, y una docena de gallinas enjauladas. Yo me coloqué en un taburete y el loco ambulante se acomodó en un montón de paja, con el legionario de pie, custodiando la salida. A aquél decidí dirigirme, con un tono vehemente y muy grave: 

“¡Muy bien, loco ambulante, muy bien sabes por qué motivo estoy hoy aquí!” 

No pareció entender nada el pobre diablo, pero al menos conseguí llamar su atención y que se incorporara: 

“¡Escúchame bien, loco! Estoy aquí porque me he dado cuenta del motivo de tu llegada. ¡Primero fueron las cuarenta gaviotas en las nonas de febrero\dots y después esto! Incluso lo absurdo tiene su significado, en un mundo gobernado por los designios de los dioses. Sabe que estoy dispuesto a perdonar tus ofensas, librándote de los cargos que se te imputan\dots ¡así que piensa bien tus respuestas! Dime, loco ambulante, contéstame a estas sencillas preguntas: ¿qué significa tu presencia? ¿Cuáles serán las penurias que habremos de soportar, y las penas que lamentar, tras tu inesperada aparición y tus monumentales desvaríos? ¿Cuál será la desgracia que sucederá al absurdo?” 

Vi que el soldado miraba a los lados, con un rostro de seriedad y estupor a partes iguales; el loco estaba quieto, en silencio y ojiplático. Yo mantuve la compostura sin embargo, sin dejar de mirarle inquisitivamente. Y mis esfuerzos dieron fruto finalmente; se levantó y se me acercó bruscamente, tanto que el legionario quiso intervenir, cosa que impedí. Me coloqué igualmente a su altura, y al verme bien el rostro me dijo: 

“¿Es que la tinta hace el mensaje? Puede que sí. ¿Y el mensaje hace la palabra? Pues\dots ¡puede que también! Pero la palabra no es el contenido de la palabra, ni la flecha del arco es la diana\dots ¡ajá—le divirtió mucho al loco llegar a esta conclusión—! Así que la tinta aparece con el problema, pero no es el problema de verdad; eso mismo pasa con un loco, incluso con uno ambulante. El loco siempre parece un problema.” 

“¡O sea que no eres tú el problema, incluso después del desorden que has causado! ¿Cuál es, entonces? ¡Contesta, diablo!” 

“¿Y si hubiera otro loco? ¡O muchos locos! ¿También ellos serían el problema? ¿Y si nadie dice la verdad, no se supone que hay muchos locos? Hay muchas formas de locura, y muchas de dar la espalda a la verdad, ¡pero no hay nada más loco que dar siempre la espalda a la verdad! ¿Dónde ha quedado la verdad? ¡Cúlpame si quieres, porque traigo conmigo un loco mensaje, pero aún quedará el contenido del mensaje\dots ajá! Y si el mensaje es verdad, y su contenido peligroso\dots ¡poco importará al respecto cuán loco sea el que hace de mensajero!” 

“Venerable cuestor—se atrevió a intervenir el soldado guardián—, está muy alterado; quizá deberíamos posponer el interrogatorio. No te imaginas lo que tememos que se acabe su estado de apaciguamiento.” 

“Será posible\dots ¡mantente en silencio!—le interpelé sin dudar (¡válgame, vaya impertinente!).” 

Por suerte la importunidad del legionario no parecía haber afectado especialmente al loco, al cual volví a dirigirme sin rodeos: 

“¿Así que traes un verdadero mensaje? ¡Ya me gustaría ver eso! ¿Cómo pueden ser la verdad tus desvaríos?” 

La verdad era que el loco ambulante estaba empezando a alterarse, pero eso no hacía sino facilitar mi tarea. En su estado de embotamiento de aquellos días no me habría servido prácticamente de nada, mientras que así me compartía sus profundidades más extravagantes e interesantes. Con la voz y las manos temblorosas, y la mirada perdida, me respondió: 

“¿Dónde está la verdad\dots y dónde su peligro? La verdad está sólo al alcance de los más sabios, brillantes, académicos, filósofos, muy difícil\dots ¡la verdad es muy difícil! Pero EL PELIGRO\dots está delante de nuestros ojos, salvo por los sabios y poderosos, que no saben verlo\dots tal vez los ignorantes tampoco sabemos verlo, ¡pero al menos no somos tan idiotas de exponernos a él!” 

En ese momento creía estar muy cerca de obtener algo provechoso: aunque parezca absurdo, estaba convencido de que las palabras del loco ambulante contenían, entre tanta absurdidad, una profecía que yo debía interpretar. ¡Tales desgracias nos habríamos ahorrado , de haber yo tenido éxito en tan arduo propósito! Le interpelé intensamente de forma que respondiera sin dilación: 

“Maldita sea, loco ambulante, ¡responde! ¿Cuál es el peligro que no ven los poderosos? ¿Cuál es el riesgo que asumen sin saberlo? ¿Qué terrible amenaza nos acecha, que terminará con nuestra tan venerada paz, con la prosperidad y la felicidad de nuestros ciudadanos? ¡Responde, miserable!” 

En este momento, noté que algo afectó profundamente al malogrado loco. En un primer momento, se quedó con la mirada perdida, totalmente desentendido de su situación en la estancia. De pronto, me miró fijamente, acercando su rostro al mío; yo no me moví de mi sitio, esperando gravemente su respuesta:  

“¡Los brillantes se harán bronce! ¡El bronce se hará piedra! ¡Y todas las piedras caerán al mar\dots Llegará el día en el que todo se hundirá!\dots ¿Y lo que no brilla? ¡Pero seréis insensatos!... ¡Eso ya está hecho de piedra!” 

Tras esta lacónica respuesta el loco ambulante quedó en silencio, sentándose en el suelo en mala postura, como si estuviera claramente agotado de su esfuerzo. Yo miraba al suelo, preguntándome si acaso, como venía sospechando, las palabras de aquel pobre desgraciado podrían significar un mal augurio. En todo caso, mi nivel de preocupación era muy elevado, más aun visto que no conseguía llegar al fondo del asunto. Lo que finalmente me sacó de mis cavilaciones fue, nuevamente, la interrupción del soberanamente inoportuno escolta:

“Mi señor cuestor, tampoco entiendo yo mucho de estos temas, pero lo que dice este pobre desgraciado tal vez tenga que ver con la última devaluación de la moneda. Entiéndeme: el señor habla de los poderosos, eso podrían ser los Césares, los emperadores—hablaba entrecortado, acelerado y nervioso, de carrerilla, así que no pude cortarlo, además de por lo absolutamente pasmado que estaba—y las monedas, ya sabes, si reduces su composición en metales preciosos, pues no lo sé\dots ¿no son menos brillantes? Quizá entonces puedan caerse al mar con más facilidad, si van en un barco\dots ¿Bueno, eso tiene sentido? A ver\dots”

“¡Esto es absolutamente impertinente y vergonzoso—le interpelé con dureza al guarda, que enmudeció arrebolado—! ¡Además de inoportuno, estúpido y desnortado: a quién se le ocurren semejantes desmanes y soberanas tonterías!”

¡Esto era ya el colmo! ¿Es que todos habían perdido el norte en esa ciudad de locos y de sinvergüenzas? Eso ya era lo último que mi paciencia estaba dispuesta a consentir en aquel lugar. El soldado, tal vez avergonzado, trato varias veces de disculparse, ora diciendo que había pasado grandes penurias, ora el ambiente con aquel loco estaba enrarecido, y a todos les había afectado. ¡Yo ya no sabía quién de verdad se había vuelto loco! A aquel impertinente guarda, y a todos los que le acompañaban en el lugar y habían considerado apropiado que él custodiara el interrogatorio, los depuse de sus puestos privilegiados, enviándolos de carceleros a las mazmorras más sombrías. 

En cuento al loco ambulante, decidí que no daba lugar a más interrogatorios. La ley era, a mi entender, clara, y el servicio prestado, habida cuenta de su declaración, debía ser recompensado. Yo creía, en primer lugar, que el hombre había hecho un sincero esfuerzo contándonos sus desvaríos, luego había cumplido su primera obligación, a saber, decir la verdad. En todo caso estaba seguro, lo cuál me hacía tener aun más clara mi decisión, de que sus palabras encerraban un intrincado significado, al fondo del cual residía un asunto de máxima gravedad, que mi deber como autoridad e investigador del Imperio era resolver. Eso pensé en ese momento, y ahora, a tenor de los acontecimientos, me reafirmo en ello. Sin embargo, si lejos estoy ahora de entender cada elemento de estos sucesos, más aun lo estaba entonces. Resolví que el loco ambulante fuera liberado, se le permitiera deambular por Atenas como condición de que fuera una vez al día a un edificio del gobierno, donde se asearía y se le daría de comer. Por suerte, el Imperio tiene el poder económico de continuar su labor civilizadora, incluso para un pobre diablo. Decidí que, en caso de que cumpliera este régimen durante un año, le serían conmutados los cargos por desorden público y sedición. 

\begin{center}
\pgfornament[width=9cm]{88}
\end{center}

Un par de días después, cenaban Julia y su tío el procónsul en el comedor de su palacio. H

\chapter{Calendas de junio}

No desapareció la preocupación de Julia por el estado de Marco Valerio; ahora bien, no sabiendo lo que había le sucedido en el muelle la noche de la farra con Lucio Junio, ni tampoco habiéndole escuchado decir más barbaridades desde aquel día, no tomó partido en el asunto, y decidió dejarlo de lado. En este último punto hay que recordar que Marco Valerio apenas hablaba con nadie aquellos días, lo cual le impedía decir ninguna locura, pero no pensarla, porque ya habrás adivinado que el pobre chico se estaba volviendo loco de remate. Tampoco compartió sus temores con su tío, ignorante de la situación, que incluso apreciaba el esfuerzo que al parecer hacía su protegido por cultivarse. 

Fue entonces que Marco Valerio se decidió a salir; iba desmarañado, mal vestido, parecía desnortado, y evitó usar la entrada principal, no queriendo ser visto por nadie. Se dirigió a la escuela del que es uno de los ciudadanos predilectos de Atenas, cuyo nombre hoy quizá pueda oírse incluso en Roma: el filósofo Cleantes. Su palabra es respetada hasta ser tenida por divina; los atenienses acuden a verlo antes de consultar el oráculo o visitar un templo, pueda resultar o no chocante y escandaloso. De hecho yo, que cada vez temo más, como ya he tenido ocasión de compartir contigo, la deriva degenerada y laxa de nuestra sociedad, he considerado a menudo a Cleantes un claro exponente, si no responsable, de ella, lo cual no ha podido sino verse reforzado con la constancia que he tenido de sus conversaciones con Marco Valerio. Juzgarás tú mismo cuánto de criminal y escandaloso hay en las mismas. 

Fue así que tal día, que debió de ser poco más de una semana antes de calendas de abril, Marco Valerio se presentó ante el filósofo Cleantes. La escuela que éste dirigía era un edificio grande aunque sencillo, quizá una antigua casa plebeya, hecho de ladrillo, y rodeado de un prado que servía de jardín y lugar de meditación. El dinero para mantener este lugar se obtenía naturalmente de donativos de mercaderes de la zona. Marco Valerio fue hasta la entrada del edificio, donde sólo hizo falta decir su nombre para que el sabio Cleantes lo recibiera. Anteriormente ya habían tenido encuentros, en los que este último quedó con la impresión de haber tratado con un joven instruido, profundo e inteligente, además de honrado y diligente. Es así que, obviamente, le extrañó sobremanera que se le presentara desarreglado y con aspecto tan agitado, lo cual no melló sin embargo en el agrado producido por su llegada, ni en la disposición a recibirlo. Le hizo subir al segundo piso de la casa, donde tenía preparada una sala de audiencias privada; ambos se sentaron frente a frente, quedando a la derecha de Marco Valerio un gran ventanal, que debía de costar una verdadera fortuna, y que dejaba a la vista toda la ciudad, hasta el puerto, iluminada por el radiante sol de primavera. Era así una estancia espaciosa y luminosa. Cleantes se dispuso a tomar la palabra: 

“Bien, hijo mío; siempre me llena de alegría el verte por aquí, aunque el aspecto que traes no invoca precisamente sentimientos placenteros\dots” 

No pudo en todo caso acabar la frase, porque de repente se abrió la puerta de golpe; entró sudoroso, sofocado y lleno de nervios un discípulo del sabio. Marco Valerio y Cleantes se dieron la vuelta para escuchar lo que tenía que decir: 

“Señor mío —balbuceaba—\dots no puedes imaginar lo que está pasando\dots el loco piojoso\dots ¡quiere entrar aquí, y no hay manera de contenerlo; está furioso\dots furioso! ¡Ese vagabundo andrajoso, que ensucia las calles y arma escándalo tras escándalo! Sabes que todos los días se le ve con un árbol del recinto de Zeus, ora dándole besos y tiernas caricias, ora discutiendo con él a grito pelado\dots ¡por tener a bien creerse que es su esposa! ¡Y lo peor es que ahora el tipo dice que quiere hablar contigo, y no se marchará sin conseguirlo; no te imaginas la escandalera que está formando! ¿Qué se supone que puede querer\dots y cómo le decimos que se vaya\dots que aquí no admitimos locos\dots? ¡Desde cuándo en la casa de la ciencia se aconseja a los chiflados!” 

Cleantes miró a su discípulo; estaba algo entrado en años, era rechoncho y calvo, y llevaba una túnica vieja pero limpia y bien cuidada: 

“Que entre —dijo—.” 

“¡¿Cómo—el discípulo regordete puso los ojos como dos rodelas de marfil—?! ¡Están todos intentado impedirlo!” 

“Ajá\dots ¡pues parece que de poco ha servido!” 

En efecto, el loco piojoso había logrado zafarse de quienes lo retenían; desde arriba se oían sus zancadas al subir las escaleras, seguidas de una ola de gritos de angustia. Al final se presentó en la mismísima sala de audiencias; apartó de un empujón al discípulo de Cleantes, y empezó a moverse de un lado a otro, tirando al suelo asientos, papeles y todo tipo de chismes que se fue encontrando. Abría la boca, enseñando sus descompuestos dientes inferiores, y tenía los ojos abiertos como platos, además de unas cejas frondosas, un rostro moreno, barbudo y mugriento; al mismo tiempo lo cubrían unos harapos ciertamente repugnantes. Toda la multitud que acudía persiguiéndolo se agolpó a la entrada de la sala; el mismo Cleantes les prohibió el paso: el mismo al que, habiéndose detenido de repente, comenzó a dirigirse el loco piojoso: 

“¡Tú\dots! ¡Ah\dots eres tú sin duda!” 

“¿Soy yo\dots?” 

“¡El hombre sabio! Los hombres son ignorantes por lo general\dots ¡y estultos hasta la médula! Pero entonces yo he preguntado: ¿en algún sitio puedo encontrar a alguien sabio\dots? Pues\dots ¿sabéis qué?\dots ¡me han dicho que venga aquí! Y ahora que te veo sé que era por ti que me han hecho venir\dots ¡porque tú tienes aspecto de saber algunas cosas, ya lo creo!” 

“Sea así; no discutiré tu juicio ni el de quienes te han hecho venir aquí; ¿cómo es que de repente vienes en busca de la sabiduría?” 

Dicho esto, Cleantes miró de reojo a Marco Valerio, para asegurarse de que no perdía la paciencia por ver su entrevista interrumpida por la aparición del loco. No sólo no era el caso, sino que incluso parecía estar disfrutando la situación, y devolvió la mirada complacido. El loco volvió entonces a dirigirse a Cleantes, pero esta vez con la mirada perdida y emitiendo extraños gruñidos mientras hablaba: 

“¡Es mi mujer\dots ella es el problema! —esto generó risas de la gente congregada junto a la estancia—. ¡Cómo no voy a quererla\dots yo la quiero\dots pero siempre está empeñada en\dots en lo mismo y lo mismo!” 

“Pérbolo, hazme un favor y cierra la puerta.” 

Su discípulo, que tal nombre tenía, hizo lo propio, empujando con brusquedad la puerta al cierre contra la fútil resistencia de los mirones de fuera. Algunos de ellos protestaron, pero sin resultado, y también sin dejar de avergonzarse por ello al de poco. Así quedaron dentro los cuatro: Marco Valerio, Cleantes, el loco piojoso y el discípulo Pérbolo. Cleantes continuó por tanto con tal prometedora entrevista: 

“Sea así: tu mujer está\dots empeñada\dots ¿y qué es aquello, si se puede saber, que tan terribles desavenencias os ha causado?” 

“¡Ah\dots!—el loco se movía sin parar de un lado a otro, hablaba entrecortado, entre dientes y con extraños cambios de volumen—. ¡Ella dice siempre que el hombre es un animal feo\dots! A veces ve pasar a gente y dice\dots ¡dice que son feos!\dots ¡según ella el humano es de los animales más feos\dots!” 

“Tú sin embargo piensas —intervino Pérbolo animado— que el hombre es un animal lleno de belleza, y de ahí el desacuerdo\dots” 

El loco se le puso delante; se le acercaba mientras hablaba, compartiéndole su fétido aliento y haciéndole retroceder espantado y, naturalmente, rebosante de creciente indignación. Sin casi poder contener la risa le dijo: 

“¡¡Estulticia máxima; el hombre no es feo; es horrible!! ¡¿Hay algo sobre la tierra que soporte tan tranquilamente la miseria?! ¡Sin honor, se arrastra por el barro!\dots Tiene que participar en su lucha\dots porque no le importa el precio a pagar\dots ¡no le importan el barro ni las mentiras, y le falta el pudor para disimularlo!\dots ¡¿Cómo no iba a ser una animal espantoso?!\dots ¡está hecho para la más infame de las batallas\dots y es horrible, por dentro y por fuera! Es un mono\dots ¡un mono!\dots ¡pero al menos los monos no pueden hablar\dots al menos no son capaces de creerse que han dejado de ser monos! ¡Al menos los monos no son capaces de ser imbéciles, a la vez que repugnantes!” 

“¿Alguna vez has visto un hombre abierto por las tripas?—preguntó, jocoso, Marco Valerio—.” 

El loco, sorprendido, se dio la vuelta de un salto: 

“¡UNA!\dots una vez vi un mono\dots ¡un MONO!\dots Y una vez\dots ¡vi a un hombre con las tripas todas abiertas! ¡¡Era espantoso y vomitivo!!” 

Se apuntó al pecho con el dedo, y empezó a deslizarlo hacia su abdomen, como si quisiera enseñarle a Marco Valerio el corte del hombre que había visto, mientras casi se cae al suelo de la risa. Esto provocó ya la absoluta indignación de Pérbolo, que se interpuso entre su maestro y el loco, reprimiendo del modo siguiente a este último: 

“¡Esta es la mayor infamia que se ha podido acercar a este santo aposento y a este santísimo y sabio hombre! ¡¡Ignominia, escándalo!!” 

Entonces se calló un momento porque había perdido la voz, y se apoyó sobre sus rodillas, jadeante, quizá arrepentido de haber perdido los nervios. Cleantes estaba en absoluto silencio, pensando profundamente. Lo que estaba sucediendo, tanto lo que había dicho el loco piojoso como la reacción de su discípulo, le impactaron e interesaron profundamente. Marco Valerio por su parte parecía estar pasándolo de perlas, mientras el loco no se dio ni por aludido durante la bronca de Pérbolo, ya que estaba mirando a todos lados y riéndoles, porque aquella silla le parecía pequeña, o aquel busto muy feo. Daba vueltas por la habitación sin mirar a nadie. Ante el silencio, y recuperado en parte, Pérbolo se incorporó, y volvió a dirigirse al loco, esta vez con un verdadero ataque de nervios: 

“¡Qué obra de arte hay mayor que el hombre; si alguien acaso pudiera reclamar tal autoría, sería el mejor poeta, el mejor escultor, el mejor músico, el mejor artista jamás sobre el mundo\dots y si cabe de todos los mundos, o fuera del mundo\dots fuera de cualquier mundo posible e imaginable! ¡Pero qué habremos de esperarnos! Un hombre que vive en la absoluta miseria: mugriento, falto de víveres y cubierto de harapos, ¡cómo iba a considerar al resto mínimamente estimable, siendo él mismo un saco de huesos repugnante! ¡Cada uno se hace un juicio de los demás según lo que es él mismo! ¡¡Un esperpento lleno de mugre viene a decirnos que la humanidad es sucia y fea!! ¡¿Quién nos dirá que estamos equivocados en todo, que no respetamos la verdad?! ¿Quién sino un loco\dots un verdadero tarado\dots un loco? ¡¡Fuera de aquí\dots LOCO!!” 

Ante este apelativo tan contundente, el loco se detuvo de repente. Se dio la vuelta y empezó a mirar a Pérbolo, y luego a los demás, con una extraña sonrisa maliciosa. Se fue acercando a ellos, deteniéndose a poca distancia y diciendo: 

“¡Dicen! Sí, dicen, que yo soy loco\dots ¡porque según ellos lo que pienso no es verdad! ¡ACÁ!\dots ¿acaso no se dice que nadie está en posesión de la verdad? ¡A mí me dijeron eso\dots una vez cuando era pequeño! Dicen que los locos cometemos errores\dots ¡pero los que no son locos los cometen más que los locos! ¡¡Máquinas de errores, fraguas de confusión, fábricas de mentiras; inundados por la mentira\dots eso es lo que hace del humano algo realmente repugnante!! Gula, pasión, interés\dots movidos por el interés mienten y se engañan, mienten y meten la pata, mienten y se cubren de vergüenza\dots ¡no les importa arrastrarse por el barro (¿Y el mito? ¿La mentira no es el mito?)! ¡¡No!! ¿Y entonces qué les pasa?\dots Les pasa que les avergüenza que alguien haya descubierto sus mentiras\dots ¡y se arrodillan para pedir clemencia! Se arrodillan para limpiarse las manos manchadas de barro\dots ¡¡con más barro!!” 

Después se arrojó al suelo y se puso a arrastrarse de modo harto esperpéntico; ¿podrá saberse por qué aquellos días la ciudad estaba llena de chiflados?  Pérbolo se tiraba de los pelos indignado, e intentó sacar a Cleantes de su reflexiva pasividad: 

“Maestro\dots ¡maestro! ¿Qué hacemos con este energúmeno? No podemos dejarlo aquí\dots” 

El loco seguía en el suelo, riéndose, cantando, haciendo todo tipo de bobadas. Cleantes, sin perder su semblante caviloso ni su mirada perdida, se levantó y le dijo a su discípulo: 

“Pérbolo, eres uno de mis más preciados aprendices\dots pero prefiero que esperes fuera: lo que voy a decir podría confundirte\dots puede malinterpretarse.” 

La cara del viejo y modesto aprendiz pareció un poema. Sin embargo, ante la insistencia de Cleantes, con el gesto y la mirada, no tuvo más remedio que marchar. Antes de irse por la puerta dijo: 

“¿Y el joven, será mejor que me lo lleve conmigo?” 

“Marco Valerio es un joven interesantísimo; quiero que se quede aquí.” 

Pérbolo se quedó mudo y salió. Sinceramente, sabiendo de este suceso se reforzó en mí la tesis de que Cleantes también estaba perdiendo el norte. ¿Cómo pudo interesarse por las estupideces y locuras del loco, y por qué demonios quería quedarse a solas con él y con Marco Valerio? Debió temer por la influencia que pudiera estar ejerciendo sobre el joven, pero actuó exactamente en el sentido contrario. Para más inri, expulsó a Pérbolo de la conversación, cuando era el único que había intervenido con un mínimo de fundamento, diciendo, como no podía ser de otra manera, que todo aquello estaba siendo un espectáculo absurdo e ignominioso. Esto, sumado a lo que dijo después en presencia de Marco Valerio, me hace pensar que había dejado de estar en sus cabales. 

Es así que Marco Valerio se quedó solo con Cleantes y con el loco, que seguía en el suelo. El filósofo se le dirigió entonces: 

“¿Opinas pues—le miró directa e inquisitivamente— que el hombre es el más feo de todos los animales, como hace este loco?” 

“Los demás animales son también horribles, pero al menos no van por ahí diciendo que son hermosísimos. Bajo mi perspectiva—decía Marco Valerio con toda naturalidad—, la cuestión radica en las conclusiones que la gente ordinaria suele sacar de este tipo de conversaciones, las raras veces que son capaces de entenderlas; el hecho de que la humana sea una especie espantosa, lamentable e hipócrita, por muchas máscaras del bien y del mal que quiera usar para ocultarlo, no implica que deba ser espantosa, ni que haga nada de bueno tolerando su indignidad.” 

“Exactamente —concluyó, tajante, Cleantes—.” 

Entonces se dirigió al loco, y le dijo con solemnidad: 

“Tú que hoy has traído la verdad a esta casa ignorante —el loco se tomó la libertad de darse por aludido, y se levantó mirando a Cleantes con los ojos abiertos como platos—, he de saber si tienes o alguna vez tuviste un nombre con el que denominarte.” 

“¡Ah —gritó el loco—\dots! ¡¿Qué más me da a mí si tengo un nombre o no?!” 

“Seas entonces sin nombre alguno, al menos tendrás un encargo; ve con tu mujer, ¡ve con ella corriendo y sin decir nada a nadie; y cuando hayas llegado dile que en vuestras discusiones eres tú el que tiene la razón de todo! ¡No somos feos; somos espantosos! ¡Mírame a mí! Piensa en mis ropas bien cuidadas, en mis palabras respetadas y tildadas de sabias, mi reputación intachable, mi generosidad y dulzura \dots ¡son un disfraz\dots sí, son todas un disfraz! ¡Debajo de ellas se oculta una realidad tan oscura, tan malvada y egoísta, llena de ambición, de confusión, de mentiras, como en los más repugnantes gorrones, ladrones y vagabundos!  

¡Siempre hablan de los buenos y los malos!\dots ¡El mundo no se divide en malos y buenos, sino en malos que triunfan y malos que fracasan! De todos modos, escúchame bien: esto no puede saberlo cualquiera, ¡entiéndeme bien! Las palabras delicadas y precisas son un peligro en manos de la masa, y si éste que soy yo es un disfraz, pues bien ¡benditos sean algunos disfraces! Es cierto que nos ocultan la maldad, pero también la alejan de nosotros, y mejor será una bondad de mentira que la maldad. Así que ve con tu mujer y disfruta tu triunfo, tú, loco lleno de sabiduría, ¡y que no se vuelva a hablar del asunto!” 

El loco, oyendo estas palabras, y entendiendo vete tú a saber qué, empezó a dar saltos de alegría, estando a punto de tirar todos los muebles de la sala. Entonces Marco Valerio, ya no divertido sino realmente excitado, se le puso en la misma cara y le dijo: 

“¿Y opinas que tu mujer también es fea?” 

El loco se quedó donde estaba, tambaleándose, como sorprendido de que alguien se le hubiera acercado tanto, y respondió: 

“¡Ajá!\dots Mi mujer también es muy fea pero\dots ¡quién se lo dice!” 

Marco Valerio sonrío maliciosamente de oreja a oreja; entonces se abrió la puerta de la estancia, y entraron todos los que esperaban al otro lado de ella, invadiendo por completo la sala. Algunos eran discípulos o estudiantes de la escuela de Cleantes, otros simplemente estaban ahí por curiosidad, ya que la noticia de que el loco había acudido al lugar se había extendido como la peste. También entró Pérbolo, indignado de que tanta gente ocupara el lugar, intentando echar a algunos, y que otros no se acercaran a tal cortina, ni a aquel busto tan preciado. Cleantes, sin embargo, no estaba alterado en absoluto y ni siquiera se levantó de su sitio. 

El loco, por su parte, no podía estar más contento de que entraran a preocuparse de lo que decía; corría de un lado a otro, acercándose y sonriendo a algunos de los que habían entrado, haciendo naturalmente que reaccionaran con terror. Hubo un momento en el que se cansó de moverse, se arrodilló en medio de la sala y, levantando los brazos hacia lo alto, gritó: 

“Mataréis dioses\dots luego os inventaréis otros que no existen, un día incluso os creeréis vosotros mismos que sois dioses\dots ¡pero siempre seréis la misma criatura, estúpida e infame!” 

Entonces se levantó y salió, zafándose de toda la marabunta que impedía su paso, y desapareció del edificio gritando, cantando y cayéndose por el suelo a cada dos pasos que daba, con más torpeza y menos recato incluso que los borrachos. Fue entonces cuando Pérbolo ya no cupo en su indignación, y se puso a gritar enajenado: 

“¡¡Ah\dots!! ¡Esto ha sido demasiado, no quiero oír a nadie decir nada más\dots! ¡¡Todo el mundo fuera de aquí ahora mismo; y vosotros los primeros!!” 

Señaló entonces a unos aprendices jóvenes, enviados por sus padres a aquella escuela, que habían entrado allí impulsados por esa absurda gracia juvenil; obviamente no tuvieron más remedio que acatar la autoridad del discípulo superior. 

Otros sin embargo se mantuvieron en la estancia, comentando y riendo la aparición del loco; fue entonces cuando el mismo Cleantes se levantó de su sitio, transformando todo aquel absurdo en una seria solemnidad. Miró a los presentes, que se callaron al instante con su mirada, y dijo: 

“Salgan todos menos el joven Marco Valerio.” 

Se sucedió un instante de silencio, en el que algunos se miraron entre sí, y después todos salieron del lugar; Cleantes se despidió afectuosamente de Pérbolo, que le correspondió en el saludo. Entonces, por fin, el filósofo y el joven se sentaron a solas como habían pretendido. Marco Valerio estaba alterado: miraba a todos lados y de cuando en cuando se reía sin motivo; la aparición del loco piojoso sin duda había hecho mella, más profunda si cabe, en su tara: 

“Pues bien, joven, espero que disculpes los imprevistos, que sin embargo habrán de considerarse provechosos, para cualquiera con la inteligencia suficiente. Heme aquí dispuesto y contento de poder escucharte.” 

Marco Valerio se inclinó para hablarle en voz baja y más de cerca: 

“Es curioso que haya venido a vernos un loco\dots porque yo he venido precisamente a consultarte acerca de la locura.” 

“Sea bienvenida entonces tu conversación —respondió en tono parecido Cleantes—, porque a veces los actos más nobles son los que se tildan de locuras.” 

“¡Ah\dots! ¿Hasta qué punto es noble, y hasta cuál loco, abandonar el mundo, haciendo lo contrario que los que lo moran esperan de uno?” 

“¡La autoridad\dots! Piensa en nosotros, los desgraciados filósofos: cuántas horas de reflexión y trabajo, cuántas dudas, cuánto sufrimiento, para decidir si algo debe ser o no verdad, generalmente para no llegar a ningún sitio; mira sin embargo a los simples y los mentecatos: si los que les rodean piensan algo, pues más vale acabar pensando lo mismo, ¡y ahí se acaba el asunto; bendita la simpleza del poder de la autoridad! Sin embargo, mi experiencia, mi ojo profundo, me han enseñado a no dar fe de lo que dicen los que me rodean, así que quizá sea más noble que loco el desdeñar lo que esperan de uno, aunque otros lo tilden de loco e infame\dots ¡qué habrá de importarle, si no se somete a la fuerza de la autoridad!” 

“¿Del mismo modo con el pensamiento que con el obrar?” 

“El pensamiento y el obrar son la misma cosa, desde dos puntos de vista distintos, al igual que la verdad y el bien. El sentimiento, que impulsa el obrar, también puede ser verdadero o falso; dos emociones distintas a menudo se contradicen, lo cual implica que al menos una de ellas tiene que ser falsa.” 

“Lo cual es cierto —apuntó Marco Valerio—, pero no deja de ser una mera aclaración conceptual.” 

“Ciertamente; y no debemos enredarnos en los conceptos. Si estás coqueteando con la nobleza y la locura, entonces tiene que tratarse de un asunto muy grave e importante, como para que nos perdamos en un absurdo galimatías de palabras.” 

“La pregunta es entonces: ¿qué es la verdad?” 

“Más bien: ¿qué es verdad? La pregunta para la que nadie tiene la respuesta\dots” 

“Pero todos tienen una respuesta\dots y yo quiero saber cuál es la tuya\dots” 

Pareció un proyectil dialéctico esta última aserción, que calló sobre Cleantes, dejándolo mudo, serio y pensativo. Se levantó de su sitio y comenzó a andar de un lado a otro, cabizbajo; no sé yo si ya se había contagiado allí mismo de la locura de Marco Valerio, o es que ésta ya hacía tiempo que se había extendido por Atenas, y que afectaba a los dos por igual. Sin parar de moverse, gesticulando, dijo: 

“Sea pues, chico\dots ¡escúchame bien! Me da igual si dedicas tu tiempo a comprender la irrisoria tragedia humana, si las flechas traicioneras del amor se han clavado en tus entrañas, si sueñas con ser el señor de la guerra; me da igual si te da alas la euforia o te acorrala el abatimiento\dots ¡aunque sea en la más escondida y silenciosa de tus profundidades, debe estar, en ti presente, la lucha eterna entre la verdad y la mentira, entre el bien y el mal; y no encontrarás a nadie fuera de estos muros en quien tan intensa y duradera haya sido esta lucha! ¡No escuches a los enemigos del cuerpo, los doctrinarios del pecado, los defensores de los justos y los compasivos; ellos no han mirado lo suficientemente lejos!” 

“¡Nunca —Marco Valerio se dio la vuelta desde su asiento para mirar hacia donde deambulaba Cleantes—\dots sólo te escucharé a ti, lo que digas ahora\dots y quedará para siempre grabado en mi pensamiento!” 

“Graba estas dos palabras: claridad y profundidad, porque la verdad no significa, ni puede significar otra cosa, que claridad en las palabras y profundidad en el sentimiento\dots ¡muy pocas veces, Marco Valerio, encontrarás a alguien que sepa cuál es el significado de las palabras\dots.! ¿Que los demás dicen esto? ¡Pues yo también digo lo mismo!\dots ¡Y así es como los superfluos usan el lenguaje! ¡La barrera de la autoridad en el lenguaje, supérala y serás el hombre más poderoso y libre de la tierra\dots decidirás lo que puede decirse y dejar de decirse, y aparecerás hasta en el pensamiento de los príncipes y los emperadores!” 

“¡¡Enséñame cómo superarla; lo quiero\dots lo quiero todo\dots necesito ante todo el poder sobre la faz de la tierra!!” 

La conversación era un concurso de atrocidades y chaladuras; Marco Valerio había alcanzado su clímax de alteración, con respecto a su ya excéntrico comportamiento de aquellos días: agitaba la cabeza en su silla, de pronto reía o se ponía serio, miraba a su interlocutor o a cualquier otro lugar de la sala indistintamente. Cleantes no quería ser menos: se movía alterado, hablaba rápido, trabado y sin vocalizar lo más mínimo, de repente gritaba en medio de una frase, o, a la vez que se quedaba clavado en el suelo, se quedaba mudo y se ponía a pensar: 

“¡Está aquí; todo está aquí al alcance de tu vista!, porque todo está fielmente conectado\dots ¡míralos! A lo que dicen los otros llaman su lenguaje, a lo que piensan verdad, a lo que ven hacer, o más bien lo que ven decir que hay que hacer, lo llaman deber\dots ¡superfluos! Sólo saben distinguir el placer del displacer\dots ¡hasta el día en que sus placeres se vuelven peligrosamente distintos a sus deberes! Es entonces que se avergüenzan, se llaman a sí mismos pecadores\dots ¡se arrepienten y se alegran de haberse arrepentido! Pues bien, ¡tú nunca te arrepientas; cada vez que culpes o te sientas culpable, estarás un paso más alejado de la verdad!\dots ¿me entiendes?” 

“¡Siempre!” 

“Sea pues, sea que tú no serás como ellos, porque tú nunca seguirás al rebaño; habrás de conocer el significado de las palabras\dots y el significado de bien es: claro en las palabras, profundo en el sentimiento\dots ¡las dos cosas, claridad y profundidad! Los buenos y los justos no son capaces de ser verdaderos, y sólo hablan de su bondad y de sus buenos sentimientos\dots ¡como si el bien pudiera ser cosa de todo el mundo! Los amantes de la investigación y la ciencia se creen que basta con la claridad\dots ¡y a eso le llaman verdad, ellos carentes de vista y de tacto, mendigos de vanidad! ¿Me entiendes lo que te digo, joven?” 

“¡Siempre y siempre!—Marco Valerio incluso se levantó al responder esto—.” 

“Da igual dónde vayas, tú mira más lejos y más profundo; nuestras pasiones nos mueven, y los simples no hacen más que seguirlas allá donde los lleven, pero tú tienes que mirar más lejos. La pasión no sólo la sientas; compréndela y entonces podrás dominarla. Los dioses son libres, pero también obran bajo el yugo de la necesidad\dots ¡domina la ciencia de las causas; conócete a ti mismo como una causa y una consecuencia, como un eslabón de la cadena eterna; conócete desde el punto de vista de la eternidad! Así habrás de dominarte a ti mismo, y sólo así puede uno aspirar a dominar el mundo entero\dots ¡sea!” 

Se quedaron los dos en silencio, con Cleantes andando bruscamente de un lado a otro, casi furioso, de forma que pareció que iba a ponerse a arrojar los muebles por la ventana. Marco Valerio se inclinaba sobre sí mismo, con la mirada perdida en el suelo y los ojos abiertos como platos, dándose golpes de ansiedad, en la pierna y en la barbilla. Se levantó de repente y preguntó: 

“¡¿Y acaso se ve compasión, cuando se mira tan lejos y tan profundo?!” 

“La justa y necesaria —dijo Cleantes casi sin pensar—; los hombres podrán ser como niños, podrán tener los caprichos de los niños\dots ¡pero uno tampoco debe ir por el mundo ignorando y atormentando a los niños! ¿Acaso se debe, sin embargo, ser indulgente con los niños malvados y maleducados, acaso no se deben corregir sus errores mediante el castigo, sea o no duro, ya que es necesario? Verás que hay felicidades que valen más que otras; conocerás el motivo del equivocado, y lo conocerás en su falsedad, ¡y por tanto no has de tener piedad de él de ninguna clase! ¡Que no te tiemble el pulso cuando tengas que recurrir a la crueldad! He aquí el mayor problema, el más grande error de planteamiento de todos los tiempos; los simples sacrifican la verdad por la felicidad. ¿Que el mundo les resulta un lugar injusto y espantoso, que sus vidas están llenas de miseria y vacías del menor sabor que llevarse a la boca? Entonces mirarán para otro lado, porque sufren cada vez que mantienen la mirada. ¿Que la bebida les distrae de su miseria? Entonces se lanzarán a ella, en vez de afrontar sus problemas con la mínima autenticidad. Pero tú harás lo contrario, sacrificarás, cuando se preste, tu felicidad a la verdad, porque ésta es la única, de todas las diosas, que es digna de ser honrada con sacrificios; poco valor tendrá tu felicidad, ni tu sufrimiento importancia\dots ¡felicidad llamarás a tu camino a la gloria, que estará lleno de aspereza, de dudas, de lucha, de contratiempos, del rechazo del prójimo, pero que llamarás feliz y acogerás con alegría, y entonces nada en el mundo podrá detenerte! ¡¿Qué quieres entonces hacer algo noble y loco?!\dots ¡Pues corre a hacerlo, y que no te detengan las olas del mar ni las legiones de Roma ni las hordas de Tartaria!” 

Desbordado de su demente exaltación se levantó Marco Valerio, y anduvo hasta la posición de Cleantes, que de tanto dar vueltas por la sala casi se había ido al extremo más alejado. Estaba agotado y se detuvo al acercarse Marco Valerio: 

“Entonces —le dijo éste—\dots ¿es todo cuestión de mirar profundo?” 

“Y de hablar claro\dots” 

“Y si miro tan profundo\dots ¿que dejo de ver nada?” 

“Quizá puedas alegrarte entonces.” 

Entonces Marco Valerio le estrechó las manos a Cleantes, se le despidió y le dijo: 

“Ahora me voy, porque aún me queda alguna cosa que pensar\dots” 

Y es así que se volvió a su casa, sin cruzar ya palabra con nadie. El comportamiento del chico siempre me resultó demasiado misterioso como para poder atribuir a esta conversación puntual con el filósofo una gran importancia al respecto de lo que acabó sucediendo en las semanas posteriores. Sin embargo, es imposible no considerar esta conversación una chaladura, ni el comportamiento de Cleantes peligroso, a la par que irresponsable. Es así que, si bien no pude tomar mayores medidas contra él, a causa de su buena reputación entre los ciudadanos de Atenas, ordené, cuando supe de todo este asunto, que nunca más volviera a entrevistarse con nadie de cuna romana. 



\chapter*{Notas}


Para más inri, esa sobrino-segunda, o lo que diantres quieras que sea, es la sobrinita mimada del mercader, así que resulta que, como nosotros, reconozcámoslo, es una fiera indómita, una nave sin ancla\dots vamos, que hace lo que quiere. Y una de las cosas que a esa chiquilla le divierte es alejarse de su buena sociedad para mezclarse con la peor calaña (exactamente como nosotros, curiosamente) \dots ¡y pasa que entre esa calaña con la que se mezcla están Albina, Nepos y la maldita prima segunda! 

¡Pero la cosa es que tu flor de alabastro, si acaso esa expresión existe, está siempre con ella, y por eso está metida en el ajo hasta el fondo!

Fue entonces que Marco Valerio alcanzó el clímax de su furia, y se lanzó contra Lucio Junio prendiéndolo por el cuello: 

“¡Cerdo sinvergüenza! ¡Canalla! ¡A ella\dots a ella sí que no se te ocurrirá mezclarla con semejante gentuza! ¡No lo permitiré!” 

Lucio Junio casi se cae al suelo, pero no por Marco Valerio sino por el ataque de risa que éste le había producido: 

“¡Es increíble; eres tan tonto como pensaba! ¡Estás a los pies de esa chica; ella es una esclava, y tú eres tan tonto que quieres ser esclavo suyo también! Podrías agradecer que se mezcle con esa gente\dots ¡para qué te crees que un tipo como el mercader Régulo habría querido utilizar a una esclava joven y hermosa!” 

“¡No te atreverás, sucio borracho; no te atreverás!” 

Y entonces Marco Valerio tiró a Lucio Junio al suelo, poniéndose encima de él. Si bien habría pensado cualquiera de frente a la escena que iban a darse una paliza, era tal el ataque de risa que dominaba a Lucio Junio que a Marco Valerio se le cortó en seco la idea de pegarle. ¿No crees que sería poco elegante aporrear a alguien que casi no puede respirar por la risa? Lucio Junio se quedó entonces sin aliento en el suelo, con su amigo dando vueltas por la explanada, pegando patadas al suelo y maldiciendo en nombre del Caos y de todos los dioses olímpicos. 


De hecho Marco Valerio se comportó aquellos días de manera extraña; apenas hablaba con nadie, salvo acaso con Julia y el procónsul, en intervenciones lacónicas y triviales, y se encerraba horas y horas a solas. Leía tragedias griegas, filosofía, fábulas egipcias y todo lo que acabara en sus manos. Tenía que empezar a conocer la vida, decía, y también que, si bien en los libros no estaba todo lo que había que aprender en una vida, cómo iba uno a entender la vida sin entender primero un mísero libro. 

Compadécete, soberano, puesto que me precio de ser tu suplicante.
Así dijo; hizo este cesar al punto su corriente, retirando las olas, e hizo la calma delante de él, llevándolo salvo a la misma desembocadura.

Entretanto la bien fabricada nave llegó velozmente a la isla de las dos Sirenas —pues la impulsaba próspero viento—. Pero enseguida cesó éste y se hizo una bonanza apacible, pues un dios había calmado el oleaje.

Levantáronse mis compañeros para plegar las velas y las
pusieron sobre la cóncava nave y, sentándose al remo,
blanqueaban el agua con los pulimentados remos.

\end{document}